\section{Proof of high dimensional recovery} \label{sec-proofs-high-d}

The proof of Theorem \ref{thm-complete} can be separated into two parts:  a recovery guarantee under a set of deterministic conditions, and a proof that the random model meets these conditions with high probability.  These sufficient deterministic conditions, roughly speaking, are (1) that the graph is connected and the nodes have tightly controlled degrees; (2) that the angles between pairs of locations is uniformly bounded away from $0$ and $\pi$; (3) that all pairwise distances are within a constant factor of each other; (4) that there are not too many corruptions affecting any single location; and (5) that the locations are `well-distributed' relative to each other in a sense we will make precise.   Theorem \ref{thm:mainthm} in Section \ref{sec-deterministic-theorem} states these deterministic conditions formally.

We will prove the deterministic recovery theorem directly, using several geometric properties concerning how deformations of a set of points induce rotations.  
%Note that an infinitesimal rigid rotation of two points $\{t_i, t_j\}$ about their midpoint to ${\{t_i + h_i, t_j + h_j\}}$ is such that $h_i - h_j$ is orthogonal to $t_{ij} = t_i - t_j$.  We will abuse terminology and say that $\|P_{\tij^\perp} (h_j - h_i)\|$ is a measure of the rotation of a finite deformation $\{h_i, h_j\}$. Morally speaking, because ShapeFit attempts to minimizes the total rotations within a feasible deformation, no deformation of the true locations results in an equal or lower objective.  These geometric properties are as follows: 
%\red{[Would like to rearrange the whole paragraph as follows:]} 
Note that an infinitesimal rigid rotation of two points $\{t_i, t_j\}$ about their midpoint to ${\{t_i + h_i, t_j + h_j\}}$ is such that $h_i - h_j$ is orthogonal to $t_{ij} = t_i - t_j$.  We will abuse terminology and say that $\|P_{\tij^\perp} (h_i - h_j)\|$ is a measure of the rotation in a finite deformation $\{h_i, h_j\}$, and we say that $\langle h_i - h_j, \ti - \tj \rangle$ is the amount of stretching in that deformation. Using this terminology, the geometric properties we establish are:

\begin{itemize}
\item If a deformation stretches two adjacent sides of a triangle at different rates, then that induces a rotation in some edge %the rest 
of the triangle (Lemma \ref{lem:rigidity1}).  
\item  If a deformation stretches two nonadjacent sides of a tetrahedron at different rates, then that induces a rotation in some edge %the rest 
of the tetrahedron (Lemma \ref{lem:rigidity2}).
\item If a deformation rotates one edge shared by many triangles, then it induces a rotation over many of those triangles, provided the opposite points of those triangles are `well-distributed' (Lemma \ref{lem:triangles}).
\item A deformation that rotates bad edges, must also rotate good edges (Lemma \ref{lem:badgood}).
\item For any deformation, some fraction of the sum of all rotations must  affect the good edges (Lemma \ref{lem:transference}).
\end{itemize}
By using these geometric properties, we show that all nonzero feasible deformations induce a
large amount of total rotation. Since some fraction of the total rotation must be on the good edges,  the objective must increase. 

In Section \ref{sec-deterministic-theorem}, we present the deterministic recovery theorem.  In Section \ref{sec-unbalanced-motions}, we present and prove Lemmas \ref{lem:rigidity1}--\ref{lem:rigidity2}.  In Section \ref{sec-triangles-inequality}, we present and prove Lemmas \ref{lem:triangles}--\ref{lem:transference}.  In Section \ref{sec-proof-mainthm}, we prove the deterministic recovery theorem.  In Section \ref{sec-gaussians-well-distributed}, we prove that Gaussians satisfy several properties, including well-distributedness, with high probability.  In Section \ref{sec-random-graph}, we prove that \erdosrenyi graphs are connected and have controlled degrees and codegrees with high probability.  Finally, in Section \ref{sec-proof-random-theorem}, we prove Theorem \ref{thm-complete}.




%A \emph{shape} is an equivalence class of a set of $n$ vectors $T=\{t_{i}\}_{i\in[n]}$
%in $\mathbb{R}^{d}$ where two sets $\{t_{i}\}_{i\in[n]}$ and $\{t_{i}'\}_{i\in[n]}$
%are equivalent if and only if $t_{i}'=\alpha(t_{i}+\vec{c})$ for
%some scalar $\alpha\in\mathbb{R}$ and vector $\vec{c}\in\mathbb{R}^{d}$.
%For a set $T$, define $L(T)=\sum_{ij}\langle t_{i}-t_{j},v_{ij}\rangle$
%and $R(T)=\sum_{ij}\|\proj_{v_{ij}^{\perp}}(t_{i}-t_{j})\|$.
%We consider the following optimization problem. 
%
%\begin{quote}
%Minimize $R(T)$ subject to $L(T)=1, \quad \sum_{i=1}^n t_i = 0$.
%\end{quote}

%%
%%
%%
%%
%%
%% Deterministic recovery condition
%%
%%
%%
%%
%%
%%


\subsection{Deterministic recovery theorem in high dimensions} \label{sec-deterministic-theorem}

To state the deterministic recovery theorem, we need two definitions.

\begin{definition}
We say that a graph $G([n],E)$ is \emph{$p$-typical} if it satisfies the following
properties:
\begin{enumerate}
  \setlength{\itemsep}{1pt} \setlength{\parskip}{0pt}
  \setlength{\parsep}{0pt}
	\item $G$ is connected,
	\item each vertex has degree between $\frac{1}{2}np$ and $2np$, and
	\item each pair of vertices has codegree between $\frac{1}{2}np^{2}$ and
$2np^{2}$, where the codegree of a pair of vertices $i,j$ is defined as $|\{k \in [n]\,:\, ik,jk \in E(G)\}|$.
\end{enumerate}
\end{definition}
Note that if $G$ is $p$-typical, then its number of edges is between $\frac{1}{4}n^{2}p$ and $n^{2}p$.


\begin{definition}
Let $T=\{t_{i}\}_{i\in[n]}\subseteq\mathbb{R}^{d}$ be a set of $n$
vectors. Let $G$ be a graph with vertex set $[n]$.
\begin{itemize}
  \setlength{\itemsep}{1pt} \setlength{\parskip}{0pt}
  \setlength{\parsep}{0pt}
\item[(i)] For a pair of vectors $x,y\in\mathbb{R}^{d}$ and a positive real
number $c$, we say that $T$ is \emph{c-well-distributed with respect
to $(x,y)$} if the following holds for all $h\in\mathbb{R}^{d}$:
\[
	\sum_{t \in T}\|\proj_{\Span\{t-x,t-y\}^{\perp}}(h)\|_2 \ge c|T|\cdot\|\proj_{(x-y)^{\perp}}(h)\|_2.
\]
\item[(ii)] We say that $T$ is \emph{$c$-well-distributed along $G$} if
for all distinct $i,j\in[n]$, the set $S_{ij}=\{t_{k}\,:\, ik,jk\in E(G)\}$
is $c$-well-distributed with respect to $(t_{i},t_{j})$.
\end{itemize}
\end{definition}


We now state sufficient deterministic recovery conditions on the graph $G$, the subgraph $\Eb$ corresponding to corrupted observations, and the locations $\Tnot$.

\begin{theorem} \label{thm:mainthm}
Suppose $\Tnot, \Eb, G$ satisfy the conditions
\begin{enumerate}
  \setlength{\itemsep}{1pt} \setlength{\parskip}{0pt}
  \setlength{\parsep}{0pt}
	\item The underlying graph $G$ is $p$-typical,
	\item For all distinct $i,j,k \in [n]$, we have $\sqrt{1-\langle\hat{t}_{ij}^{(0)},\hat{t}_{ik}^{(0)}\rangle^{2}}\ge\beta$,
	\item For all $i,j,k,\ell$ with $i\neq j$ and $k\neq l$, we have $c_{0}\|t_{k\ell}^{(0)}\|_2\le\|t_{ij}^{(0)}\|_2$,
	\item Each vertex has at most $\varepsilon n$ edges in $\Eb$ incident to it,
	\item The set $\{t_{i}^{(0)}\}_{i \in [n]}$ is $c_{1}$-well-distributed
along $G$,
	\item All vectors $\toi$ are distinct,
\end{enumerate}
for constants $0< p, \beta, c_{0}, \eps, c_{1} \leq 1$.  If $\eps \leq \frac{\beta c_0 c_1^2 p^4}{3\cdot 256 \cdot 64 \cdot 32}$, then $L(\Tnot)\neq 0$ and $\Tnot / L(\Tnot)$ is the unique optimizer of ShapeFit.
\end{theorem}


Note that Condition 3 implies that for $\muinf = \max_{i\neq j}\|t_{ij}^{(0)}\|_2$, we have $c_{0}\muinf\le\|t_{ij}^{(0)}\|_2\le\muinf$ for all distinct $i,j \in [n]$. Also note that Conditions 1--6 are invariant under translation and non-zero scalings of $\Tnot$.

Before we prove the theorem, we establish that $L(\Tnot)\neq 0$ when $\eps$ is small enough.  This property guarantees that some scaling of $\Tnot$ is feasible and occurs, roughly speaking, when $|\Eb| < |\Eg|$.


\begin{lemma} \label{lem:T0-feasible} 
%\red{If $|E_b| < c_0 |\Eg|$, then $L(\Tnot) \neq 0$.  Further, if $\varepsilon < \frac{c_0 p}{8}$, then $|E_b| < c_0 |\Eg|$.}
If $\varepsilon < \frac{c_0 p}{8}$, then $L(\Tnot) \neq 0$.
\end{lemma}
\begin{proof}
Since $v_{ij} = t_{ij}^{(0)}$ for all $ij \in E_g$, we have
\[
L(T) = \sum_{ij \in E(G)} \langle \toij, v_{ij} \rangle \geq \sum_{ij \in E_g} \| \toij\|_2 - \sum_{ij \in E_b} \| \toij \|_2.
\]
By Condition 3,  $c_0 \muinf |E_g| \leq \sum_{ij \in E_g} \| \toij\|_2$ and $\muinf |E_b| \geq  \sum_{ij\in E_b} \| \toij\|_2$.  Thus it suffices to prove that $c_0 |E_g| > |E_b|$. 
As $\varepsilon < \frac{p}{8}$, Condition 1 and 4 gives $|E_g| \geq \frac{1}{4}n^2 p - \varepsilon n^2 \geq \frac{1}{8}n^2 p$. Since $|E_b| \leq \varepsilon n^2$, if $\varepsilon < \frac{c_0 p}{8}$, then we have $c_0 |E_g| > |E_b|$. 
\end{proof}

The proof of Theorem \ref{thm:mainthm} appears in Section \ref{sec-proof-mainthm}.

%%
%%
%%
%%	SECTION : Rotational motions
%%
%%
%%
%\section{Rotational motions}

\subsection{Unbalanced parallel motions induce rotation} \label{sec-unbalanced-motions}


\begin{lemma} \label{lem:rigidity1}
Let $d\ge2$. Let $t_{1},t_{2},t_{3} \in \R^d$ be distinct.  Let $v_{1},v_{2},v_{3}\in\mathbb{R}^{d}$ and $\alpha \in \R$.   Let $\{\deltatilde_{ij}\}$ be such that  ${\langle v_{i}-v_{j}-\alpha t_{ij},\hat{t}_{ij}\rangle=}$ $\deltatilde_{ij}\|t_{ij}\|_2$ for each distinct $i,j\in[3]$. Then
\[
	\sum_{\substack{i,j\in[3]\\i<j}} 
	\|\proj_{t_{ij}^{\perp}}(v_{i}-v_{j})\|_2
	\ge \sqrt{1-\langle \hat{t}_{12},\hat{t}_{23}\rangle^{2}}\|t_{12}\|_2\left|\deltatilde_{12}-\deltatilde_{13}\right|.
\]
\end{lemma}
\begin{proof}
Note that $t_{ij}=-t_{ji}$ and $\deltatilde_{ij}=\deltatilde_{ji}$ for each
distinct $i,j\in[3]$. Define $W=\Span(\hat{t}_{12},\hat{t}_{23},\hat{t}_{31})$
and define $w_{i}=\proj_{W}v_{i}$ for each $i$. Note that
\[
\sum_{i<j}\|\proj_{t_{ij}^{\perp}}(v_{i}-v_{j})\|_2\ge\sum_{i<j}\|\proj_{t_{ij}^{\perp}}(w_{i}-w_{j})\|_2.
\]
The given condition implies $\proj_{t_{ij}^{\perp}}(w_{i}-w_{j})=w_{i}-w_{j}-(\alpha+\deltatilde_{ij})t_{ij}$
for each distinct $i,j \in [3]$. Therefore,
\begin{eqnarray*}
\sum_{i<j}\|\proj_{t_{ij}^{\perp}}(w_{i}-w_{j})\|_2 & = & \sum_{i<j}\left\| w_{i}-w_{j}-\left(\alpha+\deltatilde_{ij}\right)t_{ij}\right\|_2 \\
 & \ge & \left\| \sum_{(i,j)=(1,2),(2,3),(3,1)}w_{i}-w_{j}-\left(\alpha+\deltatilde_{ij}\right)t_{ij}\right\|_2 \\
 & = & \|\deltatilde_{12}t_{12}+\deltatilde_{23}t_{23}+\deltatilde_{31}t_{31}\|_2.
\end{eqnarray*}
Since $\deltatilde_{13}(t_{12}+t_{23}+t_{31})=0$, the right-hand-side
above equals $\|(\deltatilde_{12}-\deltatilde_{13})t_{12}+(\deltatilde_{23}-\deltatilde_{13})t_{23}\|_2$.
Furthermore,
\begin{eqnarray*}
\left\| (\deltatilde_{12}-\deltatilde_{13})t_{12}+(\deltatilde_{23}-\deltatilde_{13})t_{23}\right\|_2  & \ge & \min_{s\in\mathbb{R}}\|(\deltatilde_{12}-\deltatilde_{13})t_{12}-st_{23}\|_2\\
 & = & \left\| \proj_{t_{23}^{\perp}}(\deltatilde_{12}-\deltatilde_{13})t_{12}\right\|_2 \\
 & \ge & \left|\deltatilde_{12}-\deltatilde_{13}\right|\|t_{12}\|_2\sqrt{1-\langle 
 \hat{t}_{12},\hat{t}_{23}\rangle^{2}}.\qedhere
\end{eqnarray*}
\end{proof}


The previous lemma is applicable only when two
disproportionally scaled edges are incident to each other. The following
lemma shows how to apply the lemma above to the case when we have
two vertex-disjoint edges that are disproportionally scaled.
\begin{lemma} \label{lem:rigidity2}
Let $d\ge2$. Let $t_{1},t_{2},t_{3},t_{4} \in \R^d$ be distinct.  Let $v_{1},v_{2},v_{3},v_{4}\in\mathbb{R}^{d}$ and  $\alpha \in \R$.  Let $\{\deltatilde_{ij}\}$ be such that $\langle v_{i}-v_{j}-\alpha t_{ij},\hat{t}_{ij}\rangle=\deltatilde_{ij}\|t_{ij}\|_2$
for each distinct $i,j\in[4]$. 
Define $\beta=\min\sqrt{1-\langle\hat{t}_{ij},\hat{t}_{ik}\rangle^{2}}$
where the minimum is taken over all distinct $i,j,k\in[4]$ except for the cases when $\{j,k\}=\{1,2\}$. Then
\[
	\sum_{\substack{i,j\in[4]\\i<j}} \|\proj_{t_{ij}^{\perp}}(v_{i}-v_{j})\|_2
	\ge \frac{\beta}{4}\|t_{12}\|_2\left|\deltatilde_{12}-\deltatilde_{34}\right|.
\]
\end{lemma}
\begin{proof}
Note that $t_{ij}=-t_{ji}$ and $\deltatilde_{ij}=\deltatilde_{ji}$ for each
distinct $i,j\in[4]$. Since the given conditions are symmetric under
re-labelling of ($1$ and $2$), and of ($3$ and $4$), we may re-label
if necessary and assume that $\|t_{13}\|_2 \ge\max\{\|t_{14}\|_2 ,\|t_{23}\|_2 ,\|t_{24}\|_2 \}$.
By the triangle inequality, we have $2\|t_{13}\|_2 \ge\|t_{13}\|_2 +\|t_{23}\|_2 \ge\|t_{12}\|_2 $.
Apply Lemma \ref{lem:rigidity1} to the triangle $\{1,2,3\}$ to obtain
\begin{eqnarray}
\sum_{i<j,\ i,j \in \{1, 2, 3\}}\|\proj_{t_{ij}^{\perp}}(v_{i}-v_{j})\|_2 & \ge & \sqrt{1-\langle \hat{t}_{12},\hat{t}_{23}\rangle^{2}}\|t_{12}\|_2 \left|\deltatilde_{12}-\deltatilde_{13}\right|\nonumber \\
 & \ge & \beta\|t_{12}\|_2 \left|\deltatilde_{12}-\deltatilde_{13}\right|,\label{eq:k4_1}
\end{eqnarray}
and similarly apply the lemma to the triangle $\{3,1,4\}$ to obtain
\begin{eqnarray}
\sum_{i<j, \ i,j\in\{1, 3, 4\}}\|\proj_{t_{ij}^{\perp}}(v_{i}-v_{j})\|_2 & \ge & \sqrt{1-\langle \hat{t}_{13},\hat{t}_{14}\rangle^{2}}\|t_{13}\|_2 \left|\deltatilde_{13}-\deltatilde_{34}\right|\nonumber \\
 & \ge & \beta\|t_{13}\|_2 |\deltatilde_{13}-\deltatilde_{34}|\ge\frac{\beta}{2}\|t_{12}\|_2 |\deltatilde_{13}-\deltatilde_{34}|.\label{eq:k4_2}
\end{eqnarray}
By adding \eqref{eq:k4_1} and \eqref{eq:k4_2}, we see that
\begin{align*}
	&\,\sum_{i<j,\ i,j \in \{1, 2, 3\}}\|\proj_{t_{ij}^{\perp}}(v_{i}-v_{j})\|_2
	+
	\sum_{i<j, \ i,j\in\{1, 3, 4\}}\|\proj_{t_{ij}^{\perp}}(v_{i}-v_{j})\|_2 \\
	&\,\ge \,\, 
	\beta\|t_{12}\|_2 \left|\deltatilde_{12}-\deltatilde_{13}\right|
	+
	\frac{\beta}{2}\|t_{12}\|_2 \left|\deltatilde_{13}-\deltatilde_{34}\right|
	\ge
	\frac{\beta}{2}\|t_{12}\|_2 \left|\deltatilde_{12}-\deltatilde_{34}\right|.
\end{align*}
The lemma follows since the left-hand-side is bounded from
above by $2 \sum_{\substack{i,j\in[4]\\i<j}} \|\proj_{t_{ij}^{\perp}}(v_{i}-v_{j})\|_2$.
\end{proof}






%For example, in $\mathbb{R}^{3}$ the well-distributedness of vectors
%$t_{1},\cdots,t_{k}$ with respect to a pair $x,y$ is determined
%by the vectors $v_{i}$ defined as a unit vector orthogonal to both
%$x-t_{i}$ and $y-t_{i}$. Since all vectors $v_{i}$ are orthogonal
%to $x-y$, we can think of $v_{i}$ as vectors on a circle in $\mathbb{R}^{2}$.
%Heuristically, the vectors $t_{1},\cdots,t_{k}$ are well-distributed
%with respect to $x,y$ if $v_{1},\cdots,v_{k}$ are geometrically
%well-distributed along the circle. For example if $k=2$ and $v_{1},v_{2}$
%are orthogonal, then $t_{1},t_{2}$ will be $\frac{1}{2}$-well-disbritued
%with respect to $x,y$.

\subsection{Triangles inequality and rotation propagation} \label{sec-triangles-inequality}

\begin{lemma}[Triangles Inequality] \label{lem:triangles}
Let $d\ge3$; $x,y,t_{1},t_{2},\cdots,t_{k}\in\mathbb{R}^{d}$.
If $T=\{t_{1},\cdots,t_{k}\}$ is $c$-well-distributed
with respect to $(x,y)$, then for all vectors $h_{x},h_{y},h_{1},\cdots,h_{k}\in\mathbb{R}^{d}$
and sets $X\subseteq[k]$, we have
\[
\sum_{i\in[k]\setminus X}\|\proj_{(x-t_{i})^{\perp}}(h_{x}-h_{i})\|_2 +\|\proj_{(t_{i}-y)^{\perp}}(h_{i}-h_{y})\|_2 \ge(ck-|X|)\cdot\|\proj_{(x-y)^{\perp}}(h_{x}-h_{y})\|_2.
\]
\end{lemma}
\begin{proof}
For each $i\in[k]$, define $W_{i}=\Span\langle x-t_{i}\,,\, t_{i}-y\rangle$.
Define $P$ as the projection map to the space of vectors orthogonal
to $x-y$, and define $P_{i}$ for each $i\in[k]$  as the projection
map to $W_{i}^{\perp}$. Since $(x-t_{i})^{\perp}\supseteq W_{i}^{\perp}$
and $(t_{i}-y)^{\perp}\supseteq W_{i}^{\perp}$, it follows that 
\begin{eqnarray*}
 &  & \sum_{i\in[k]\setminus X}\|\proj_{(x-t_{i})^{\perp}}(h_{x}-h_{i})\|_2 +\|\proj_{(t_{i}-y)^{\perp}}(h_{i}-h_{y})\|_2\\
 & \ge & \sum_{i\in[k]\setminus X}\|P_{i}(h_{x}-h_{i})\|_2 +\|P_{i}(h_{i}-h_{y})\|_2 \ge\sum_{i\in[k]\setminus X}\|P_{i}(h_{x}-h_{y})\|_2.
\end{eqnarray*}
Since $t_{1},\cdots,t_{k}$ are well-distributed with respect to $(x,y)$,
we have
\begin{equation}
\sum_{i\in[k]}\|P_{i}(h_{x}-h_{y})\|_2 \ge ck\cdot\|P(h_{x}-h_{y})\|_2.\label{eq:proj_triangles}
\end{equation}
Since $\|P_{i}(h_{x}-h_{y})\|_2\le\|P(h_{x}-h_{y})\|_2$ holds for all
$i$, it follows that 
\[
\sum_{i\in[k]\setminus X}\|P_{i}(h_{x}-h_{y})\|_2 \ge(ck-|X|)\cdot\|P(h_{x}-h_{y})\|_2,
\]
proving the lemma.
\end{proof}

The proof of Theorem \ref{thm:mainthm} will rely on the following two lemmas, which state that rotational motions on some parts of the graph bound rotational motions on other parts. The following lemma relates the rotational motions on bad edges to the rotational motions on good edges.  Recall the notation $\tij = (1 + \deltaij) \toij + \etaij \sij$ where $\sij$ is a unit vector orthogonal to $\toij$ and $\etaij = \| P_{\toijperp} \tij \|_2$.



\begin{lemma} \label{lem:badgood} 
Fix $T$. If $\varepsilon_0 \leq \frac{c_1 p^2}{8}$, then $\sum_{ij\in E_{g}}\eta_{ij}\ge\frac{c_{1}p^{2}}{8\varepsilon_0}\sum_{ij\in E_{b}}\eta_{ij}$. 
\end{lemma}
\begin{proof}
For each edge $ij \in E(K_n)$, by Conditions \ref{enu:typical}, \ref{enu:bad}, \ref{enu:welldistributed}; Lemma \ref{lem:triangles} and $\varepsilon_0 \leq \frac{c_1p^2}{8}$,
we have 
\[
	\sum_{\substack{k\neq i,j\\ik,jk\in E_{g}}}(\eta_{ik}+\eta_{jk})
	\ge\left(c_{1}\cdot\frac{1}{2}np^{2}-2\varepsilon_0n\right)\cdot\eta_{ij}
	\ge\frac{c_{1}}{4}np^{2}\cdot\eta_{ij}.
\]
Therefore, if we sum the inequality above for all bad edges $ij\in E_{b}$, then
\[
	\sum_{ij\in E_{b}}\sum_{\substack{k\neq i,j\\ik,jk\in E_{g}}} (\eta_{ik}+\eta_{jk})
	\ge\frac{c_{1}}{4}np^{2}\cdot\sum_{ij\in E_{b}}\eta_{ij}.
\]
For fixed $ik\in E_{g}$, the left-hand-side may sum $\eta_{ik}$
as many times as the number of bad edges incident to the edge $ik$.
Hence by Condition 4, the left-hand-side of above is at most 
\[
	\sum_{ij\in E_{b}} \sum_{\substack{k\neq i,j\\ik,jk\in E_{g}}} (\eta_{ik}+\eta_{jk})
	\le 2\varepsilon_0 n\cdot\sum_{ij\in E_{g}}\eta_{ij}.
\]
Therefore by combining the two inequalities above, we obtain 
\[
	\sum_{ij\in E_{b}}\eta_{ij}
	\le\frac{8\varepsilon_0}{c_{1}p^{2}}\sum_{ij\in E_{g}}\eta_{ij}. %\leq \sum_{ij\in E_{g}}\eta_{ij}.
	\qedhere
\]
\end{proof}

The following lemma relates the rotational motions over the good graph $\Eg$ to rotational motions over the complete graph $K_n$.

\begin{lemma} \label{lem:transference} 
Fix $T$. If $\varepsilon_0 \leq \frac{c_1 p^2}{8}$, then $\sum_{ij\in E_{g}}\eta_{ij}\ge\frac{c_{1}p}{16}\sum_{ij\in E(K_{n})}\eta_{ij}$.
\end{lemma}
\begin{proof}
For each $ij\in E(K_{n})$, since $\{t_{i}^{(0)}\}_{i=1}^{n}$ is
$c_{1}$-well-distributed along $G$ and $G$ is $p$-typical, we
have as in the proof of Lemma~\ref{lem:badgood},
\[
	\sum_{\substack{k\neq i,j\\ik,jk\in E_{g}}}(\eta_{ik}+\eta_{jk})
	\ge \left(c_{1}\cdot\frac{1}{2}np^{2}-2\varepsilon_0 n\right) \cdot\eta_{ij} 
	\ge \frac{c_{1}}{4}np^{2}\cdot\eta_{ij}.
\]
If we sum the above over all $ij\in E(K_{n})$, we obtain 
\[
	\sum_{ij\in E(K_{n})}\sum_{\substack{k\neq i,j\\ik,jk\in E_{g}}}(\eta_{ik}+\eta_{jk})
	\ge \frac{c_{1}}{4}np^{2} \cdot\sum_{ij\in E(K_{n})}\eta_{ij}.
\]
For a fixed $ik\in E_{g}$, the left-hand-side may sum $\eta_{ik}$
as many as times as the number of edges of $G$ incident to $ik$. Therefore
since $G$ is $p$-typical, we see that 
\[
	\sum_{\substack{k\neq i,j\\ik,jk\in E_{g}}}(\eta_{ik}+\eta_{jk})
	\le 2\cdot2np\sum_{ij\in E_{g}}\eta_{ij}.
\]
By combining the two inequalities, we obtain 
\[
 	\frac{c_{1}}{4}np^{2}\sum_{ij\in E(K_{n})}\eta_{ij}
 	\le 4np\sum_{ij\in E_{g}}\eta_{ij},
\]
and thus $\frac{c_{1}p}{16}\sum_{ij\in E(K_{n})}\eta_{ij}\le\sum_{ij\in E_{g}}\eta_{ij}.$
\end{proof}






\subsection{Proof of Theorem \ref{thm:mainthm}} \label{sec-proof-mainthm}


We now prove the deterministic recovery theorem.

\begin{proof}[Proof of Theorem \ref{thm:mainthm}]

By Lemma \ref{lem:T0-feasible} and the fact that Conditions 1--6 are invariant under global translation and nonzero scaling, we can take $\Tnotbar=0$ and $L(\Tnot) = 1$ without loss of generality.  The variable $\muinf$ is to be understood accordingly.  


We will directly prove that $R(T) > R(\Tnot)$ for all $T\neq \Tnot$ such that $L(T)=1$ and $\Tbar = 0$. Consider an arbitrary feasible $T$ and recall the notation $\tij = (1 + \deltaij) \toij + \etaij \sij$ where $\sij$ is a unit vector orthogonal to $\toij$ and $\etaij = \| P_{\toijperp} \tij \|_2$. A useful lower bound for the objective $R(T)$ is given by 
\begin{eqnarray}
	R(T)
	=\sum_{ij}\|\proj_{v_{ij}^{\perp}}t_{ij}\| _2
	& = & \sum_{ij\in E_{g}}\eta_{ij}+\sum_{ij\in E_{b}}\|\proj_{v_{ij}^{\perp}}t_{ij}\|_2\nonumber \\
	& \ge & \sum_{ij\in E_{g}}\eta_{ij}+\sum_{ij\in E_{b}}\left(\|\proj_{v_{ij}^{\perp}}t_{ij}^{(0)}\|_2-|\delta_{ij}|\|t_{ij}^{(0)}\|_2-\eta_{ij}\right)\nonumber \\
	& \ge & R(\Tnot)+\sum_{ij\in E_{g}}\eta_{ij}-\sum_{ij\in E_{b}}(|\delta_{ij}|\|t_{ij}^{(0)}\|_2+\eta_{ij}).\label{eq:gain3}
\end{eqnarray}


Suppose that $\sum_{ij\in E_{b}}|\delta_{ij}|\|t_{ij}^{(0)}\|_2<\sum_{ij\in E_{b}}\eta_{ij}$.
Since Lemma \ref{lem:badgood} for $\varepsilon\le\frac{c_{1}p^2}{16}$
implies $\sum_{ij\in E_{b}}\eta_{ij}\le\frac{1}{2}\sum_{ij\in E_{g}}\eta_{ij}$,
by (\ref{eq:gain3}), we have
\begin{eqnarray*}
	R(T) & \ge & R(\Tnot)+\sum_{ij\in E_{g}}\eta_{ij}-\sum_{ij\in E_{b}}(|\delta_{ij}|\|t_{ij}^{(0)}\|_2+\eta_{ij})\\
	& > & R(\Tnot)+\sum_{ij\in E_{g}}\eta_{ij}-\sum_{ij\in E_{b}}2\eta_{ij}\ge R(\Tnot).
\end{eqnarray*}
Hence we may assume 
\begin{equation}
	\sum_{ij\in E_{b}}|\delta_{ij}|\|t_{ij}^{(0)}\|_2\ge\sum_{ij\in E_{b}}\eta_{ij}.\label{eq:badtwist}
\end{equation}

In the case $|\Eb| \neq 0$, define $\overline{\delta}=\frac{1}{|E_{b}|}\sum_{ij\in E_{b}}|\delta_{ij}|$
as the average `relative parallel motion' on the bad edges. For distinct edges $ij,k\ell\in E(K_{n})$,
if $\{i,j\}\cap\{k,\ell\} = \emptyset$, then define $\eta(ij,k\ell)=\eta_{ij}+\eta_{ik}+\eta_{i\ell}+\eta_{jk}+\eta_{j\ell}+\eta_{k\ell}$,
and if $\{i,j\}\cap\{k,\ell\}\neq\emptyset$ (without loss of generality,
assume $\ell=i$), then define $\eta(ij,k\ell)=\eta_{ij}+\eta_{ik}+\eta_{jk}$. \\

\noindent \textbf{Case 0}. $\bar{\delta} = 0$ or $|\Eb| = 0$.
\medskip

Note that $\bar{\delta}=0$ implies $\delta_{ij} =0$ for all $ij \in E_b$, which by \eqref{eq:badtwist} implies
$\eta_{ij} =0$ for all $ij \in E_b$.  Therefore by \eqref{eq:gain3}, we have
\[
	R(T) \geq R(\Tnot) + \sum_{ij \in E_g} \eta_{ij}.
\]
If $\sum_{ij \in E_g} \eta_{ij} > 0$, then we have $R(T) > R(\Tnot)$. Thus 
we may assume that $\eta_{ij} = 0$ for all $ij \in E_g$. In this case, we will show that $T = \Tnot$.

%Therefore $\eta_{ij} = 0$ for all $ij \in E(G)$, and by
By Lemma~\ref{lem:transference}, if $\varepsilon_0 \le \frac{c_1p^2}{8}$, then
$\eta_{ij}=0$ for all $ij \in E(G)$ implies that
$\eta_{ij} = 0$ for all $ij \in E(K_n)$. 
For $ij \in E_b$, since $\delta_{ij} = \eta_{ij} = 0$, it follows that $\ell_{ij} = \ell_{ij}^{(0)}$.
Since $\delta_{ij} \|\toij\|_2 = \ell_{ij} - \ell_{ij}^{(0)}$ for $ij \in E_g$, we have
\[
	0
	= \sum_{ij \in E(G)} (\ell_{ij} - \ell_{ij}^{(0)})
	= \sum_{ij \in E_b} (\ell_{ij} - \ell_{ij}^{(0)}) + \sum_{ij \in E_g} (\ell_{ij} - \ell_{ij}^{(0)})
	= \sum_{ij \in E_g} (\ell_{ij} - \ell_{ij}^{(0)})
	= \sum_{ij \in E_g} \delta_{ij} \|\toij\|_2,
\] 
where the first equality is because $L(T) = L(\Tnot) = 1$.
By Condition 6, $\|\toij\|_2 \neq 0$ for all $i \neq j$. Therefore, if $\delta_{ij} \neq 0$ for some $ij \in E_g$, then there exists $ab, cd \in E_g$ 
such that $\delta_{ab} > 0$ and $\delta_{cd} < 0$. 
By Lemma~\ref{lem:rigidity1} or \ref{lem:rigidity2} and Condition 2, this forces $\eta(ab,cd) > 0$, contradicting the fact that
$\eta_{ij} = 0$ for all $ij \in E(K_n)$. 
Therefore $\delta_{ij} = 0$ for all $ij \in E_g$, and hence
$\delta_{ij} = 0$ for all $ij \in E(G)$.  

Define $t_i = t_i^{(0)} + h_i$ for each $i \in [n]$. Because $\eta_{ij} = \delta_{ij} = 0$ for all $ij \in E(G)$, we have $h_i = h_j$ 
for all $ij \in E(G)$. Since $G$ is connected, this implies $h_i = h_j$ for all 
$i \in [n]$. Then by the constraint $\sum_{i \in [n]} t_i = \sum_{i \in [n]} t_i^{(0)} = 0$, 
we get $h_i = 0$ for all $i \in [n]$. Therefore $T = \Tnot$. \\

\noindent\textbf{Case 1}. $\bar{\delta} \neq 0$ and $\sum_{ij\in E_{g}}|\delta_{ij}|<\frac{1}{8}\overline{\delta}|E_{g}|$ and $|\Eb| \neq 0$. 
\medskip

Define $L_{b}=\{ij\in E_{b}:|\delta_{ij}|\ge\frac{1}{2}\overline{\delta}\}$.
Note that $\sum_{ij\in E_{b}\setminus L_{b}}|\delta_{ij}|<\frac{1}{2}\overline{\delta}|E_{b}|$
and therefore 
\begin{equation}
\sum_{ij\in L_{b}}|\delta_{ij}|=\sum_{ij\in E_{b}}|\delta_{ij}|-\sum_{ij\in E_{b}\setminus L_{b}}|\delta_{ij}|>\sum_{ij\in E_{b}}|\delta_{ij}|-\frac{1}{2}\overline{\delta}|E_{b}|=\frac{1}{2}\overline{\delta}|E_{b}|.\label{eq:large_delta}
\end{equation}
Define $F_{g}=\{ij\in E_{g}:|\delta_{ij}|<\frac{1}{4}\overline{\delta}\}$.
Then by the condition of Case 1,
\[
\frac{1}{8}\overline{\delta}|E_{g}|>\sum_{ij\in E_{g}}|\delta_{ij}|\ge\sum_{ij\in E_{g}\setminus F_{g}}|\delta_{ij}|\ge\frac{1}{4}\overline{\delta}|E_{g}\setminus F_{g}|,
\]
and therefore $|E_{g}\setminus F_{g}|<\frac{1}{2}|E_{g}|$, or equivalently,
$|F_{g}|>\frac{1}{2}|E_{g}|$.

For each $ij\in L_{b}$ and $k\ell\in F_{g}$, by Lemmas \ref{lem:rigidity1}, \ref{lem:rigidity2},
and Condition 3, we have $\eta(ij,k\ell)\ge \frac{\beta}{4}|\delta_{ij} - \delta_{k\ell}| \cdot\|t_{ij}\|_2 \ge
\frac{\beta}{4} \cdot \frac{1}{2}|\delta_{ij}| \cdot \|t_{ij}\|_2 \ge
\frac{\beta c_{0}\muinf}{8}|\delta_{ij}|$. 
Therefore by Condition 1, 
\begin{eqnarray*}
	\sum_{ij\in E_{b}}\sum_{k\ell\in E_{g}}\eta(ij,k\ell) 
	& \ge & \sum_{ij\in L_{b}}\sum_{k\ell\in F_{g}}\frac{\beta c_{0}\muinf}{8}|\delta_{ij}|=\sum_{ij\in L_{b}}|F_{g}|\cdot\frac{\beta c_{0}\muinf}{8}|\delta_{ij}|\\
	& >  & \sum_{ij\in L_{b}}\frac{\beta c_{0}\muinf}{16}|E_{g}||\delta_{ij}|
	\ge\frac{\beta c_{0}\muinf}{16}|E_{g}|\cdot\frac{1}{2}\overline{\delta}|E_{b}|,
\end{eqnarray*}
where the last inequality follows from (\ref{eq:large_delta}). For
each $ij\in E(K_{n})$, we would like to count how many times each
$\eta_{ij}$ appear on the left hand side. If $ij\in E_{b}$, then
there are at most ${n \choose 2}$ $K_{4}$s and $n$ $K_{3}$s containing
$ij$; hence $\eta_{ij}$ may appear at most $6{n \choose 2}+3n=3n^{2}$
times. If $ij\notin E_{b}$, then $\eta_{ij}$ appears when there is
a $K_{4}$ or a $K_{3}$ containing $ij$ and some bad edge. By Condition
4, there are at most $2\varepsilon n$ such bad $K_{3}$s. If the
bad edge in $K_{4}$ is incident to $ij$, then there are at most
$2\varepsilon n\cdot(n-3)$ such $K_{4}$s, and if the bad edge is
not incident to $ij$, then there are at most $|E_{b}|\le\varepsilon n^{2}$
such $K_{4}$. Thus $\eta_{ij}$ may appear at most $3\cdot2\varepsilon n+6\cdot(2\varepsilon n(n-3)+\varepsilon n^{2})\le18\varepsilon n^{2}$
times. Therefore
\begin{eqnarray*}
	\sum_{ij\in E_{b}}\sum_{k\ell\in E_{g}}\eta(ij,k\ell) 
	& \le & \sum_{ij\in E_{b}}3n^{2}\cdot\eta_{ij}+\sum_{ij\in E(K_{n})}18\varepsilon n^{2}\cdot\eta_{ij}.
\end{eqnarray*}
By Lemma \ref{lem:badgood}, if $\varepsilon < \frac{c_1p^2}{8}$, we have
\[
	\sum_{ij\in E_{b}}\sum_{k\ell\in E_{g}}\eta(ij,k\ell)
	\le\frac{24\varepsilon}{c_{1}p^{2}}n^{2}\sum_{ij\in E_{g}}\eta_{ij}+\sum_{ij\in E(K_{n})}18\varepsilon n^{2}\cdot\eta_{ij}
	\le\frac{42\varepsilon}{c_{1}p^{2}}n^{2}\sum_{ij\in E(K_{n})}\eta_{ij}.
\]
Hence
\[
	\frac{42\varepsilon}{c_{1}p^{2}}n^{2}\sum_{ij\in E(K_{n})}\eta_{ij}\ge\frac{\beta c_{0}\muinf}{32}|E_{g}|\cdot\overline{\delta}|E_{b}|.
\]
If $\varepsilon< \frac{p}{8}$, then $|E_{g}|\ge\frac{n^{2}p}{4}-|E_{b}|\ge\frac{n^{2}p}{8}$. 
Further, if $\varepsilon< \frac{\beta c_0 c_1^2 p^4}{32 \cdot 42 \cdot 32 \cdot 8}$, then by Condition 3,
 $\bar{\delta} \neq 0$, and $|\Eb| \neq 0$, the above implies 
\begin{align*}
	\sum_{ij\in E(K_{n})}\eta_{ij}
	\ge&\, \frac{\beta c_{0}c_{1}p^{2}}{42\cdot 32 \varepsilon n^{2}}\muinf|E_{g}|\cdot\overline{\delta}|E_{b}| 
	\ge \frac{\beta c_{0}c_{1}p^{3}}{42\cdot 32 \cdot 8} \cdot \frac{1}{\varepsilon} \cdot \muinf \overline{\delta}|E_{b}|  \\
	>&\, \frac{32}{c_{1}p}\muinf\cdot\overline{\delta}|E_{b}|
	\ge \frac{32}{c_{1}p}\sum_{ij\in E_{b}}|\delta_{ij}|\|t_{ij}^{(0)}\|_2.
\end{align*}
Lemma \ref{lem:transference} implies 
\[
	\sum_{ij\in E_{g}}\eta_{ij}
	\ge \frac{c_{1}p}{16}\sum_{ij\in E(K_{n})}\eta_{ij}
	> 2\sum_{ij\in E_{b}}|\delta_{ij}|\|t_{ij}^{(0)}\|_2.
\]
%Lemma \ref{lem:badgood} implies $\sum_{ij\in E_{g}}\eta_{ij}\ge 2\sum_{ij\in E_{b}}\eta_{ij}$ if $\varepsilon < \frac{c_1p^2}{16}$.
Therefore by (\ref{eq:badtwist}),we have $\sum_{ij\in E_{g}}\eta_{ij}>\sum_{ij\in E_{b}}(|\delta_{ij}|\|t_{ij}^{(0)}\|_2+\eta_{ij})$ if $\varepsilon \le \min\{\frac{c_1p^2}{8}, \frac{p}{8}, \frac{\beta c_0 c_1^2 p^4}{32 \cdot 42 \cdot 32 \cdot 8} \}$.
By (\ref{eq:gain3}), this shows $R(T)>R(\Tnot)$. This condition on $\eps$ is satisfied under the assumption $\eps \leq \frac{\beta c_0 c_1^2 p^4}{3\cdot 256 \cdot 64 \cdot 32}$.
\\

\noindent\textbf{Case 2}. $\bar{\delta} \neq 0$ and $\sum_{ij\in E_{g}}|\delta_{ij}|\ge\frac{1}{8}\overline{\delta}|E_{g}|$  and $|\Eb| \neq 0$. 
\medskip

Define $E_{+}=\{ij\in E_{g}\,:\,\delta_{ij}\ge0\}$ and $E_{-}=\{ij\in E_{g}\,:\,\delta_{ij}<0\}$.
Since $\ell_{ij}-\ell_{ij}^{(0)}=\delta_{ij}\|t_{ij}^{(0)}\|_2$
for $ij\in E_{g}$, we have 
\begin{eqnarray*}
0=\sum_{ij \in E(G)}(\ell_{ij}-\ell_{ij}^{(0)}) & = & \sum_{ij\in E_{b}}(\ell_{ij}-\ell_{ij}^{(0)})+\sum_{ij\in E_{g}}\delta_{ij}\|t_{ij}^{(0)}\|_2.
\end{eqnarray*}
where the first equality follows from $L(T) = L(\Tnot)$.  Therefore, 
\begin{eqnarray*}
	\left| \sum_{ij\in E_{g}}\delta_{ij}\|t_{ij}^{(0)}\|_2 \right|
	\le \left| \sum_{ij\in E_{b}}(\ell_{ij}-\ell_{ij}^{(0)}) \right|
	\le \sum_{ij\in E_{b}}(|\delta_{ij}|\|t_{ij}^{(0)}\|_2+\eta_{ij}) 
	\le 2\muinf\overline{\delta}|E_{b}|,
\end{eqnarray*}
where the last inequality follows from (\ref{eq:badtwist}), Condition 3, and the definition of $\overline{\delta}$. On the other hand, the condition of Case 2 and Condition 3 implies 
$\sum_{ij\in E_{g}}|\delta_{ij}|\|t_{ij}^{(0)}\|_2 \ge \frac{1}{8}c_0 \muinf\overline{\delta}|E_g|$. 
Therefore
\[
	\sum_{ij \in E_-} (-\delta_{ij}) \|t_{ij}^{(0)}\|_2
	= \frac{1}{2}\left( -\sum_{ij\in E_{g}}\delta_{ij}\|t_{ij}^{(0)}\|_2 + \sum_{ij\in E_{g}}|\delta_{ij}|\|t_{ij}^{(0)}\|_2\right)
	\ge \frac{1}{2} \left(\frac{1}{8} c_0 \muinf\overline{\delta}|E_g| - 2\muinf\overline{\delta}|E_{b}|\right).
\]
If $\varepsilon \le \frac{1}{256}c_0p$, then 
since $|E_b| \le \varepsilon n^2$ and $|E_g| \ge \frac{1}{4}n^2p - |E_b| \ge \frac{1}{8}n^2p$, 
we see that $\frac{1}{8} c_0 \muinf\overline{\delta}|E_g| - 2\muinf\overline{\delta}|E_{b}| \ge \frac{1}{16}c_0\muinf \bar{\delta}|E_g|$.
Therefore $\sum_{ij\in E_{-}}(-\delta_{ij})\|t_{ij}^{(0)}\|_2 \ge\frac{1}{32}c_{0}\muinf \overline{\delta}|E_{g}|$.
Similarly, $\sum_{ij\in E_{+}} \delta_{ij}\|t_{ij}^{(0)}\|_2 \ge \frac{1}{32}c_{0}\muinf \overline{\delta}|E_{g}|$.

If $|E_+| \ge \frac{1}{2}|E_g|$, then 
by Lemmas \ref{lem:rigidity1}, \ref{lem:rigidity2}, and Condition 3, we have 
\begin{eqnarray*}
	\sum_{ij\in E_{-}}\sum_{k\ell\in E_{+}}\eta(ij,k\ell) 
	& \ge & \sum_{ij\in E_{-}}\sum_{k\ell\in E_{+}}\frac{\beta}{4}(-\delta_{ij})\|t_{ij}^{(0)}\|_2\\
	& \ge & \sum_{ij\in E_{-}}(-\delta_{ij})\|t_{ij}^{(0)}\|_2\cdot \frac{\beta}{4}|E_{+}|
	\ge \frac{\beta}{4}|E_{+}|\cdot\frac{1}{32}c_{0}\muinf \overline{\delta}|E_{g}| \\
	& \ge & \frac{\beta}{256} c_{0} \muinf\overline{\delta}|E_{g}|^2. 
\end{eqnarray*}
Similarly, if $|E_-| \ge \frac{1}{2}|E_g|$, then we can switch the order of summation and
consider $\sum_{ij \in E_+} \sum_{k\ell \in E_-} \eta(ij, k\ell)$ to obtain the 
same conclusion.

Since each edge is contained in at most $\frac{n(n-1)}{2}$
copies of $K_{4}$ and $n$ copies of $K_{3}$ (and there are 6 edges
in a $K_{4}$, 3 edges in a $K_3$), we have 
\[
	\sum_{ij\in E_{-}}\sum_{k\ell\in E_{+}}\eta(ij,k\ell)
	\le\left(6\frac{n(n-1)}{2}+3n\right)\sum_{ij\in E(K_{n})}\eta_{ij}\le3n^{2}\sum_{ij\in E(K_{n})}\eta_{ij}.
\]
If $\varepsilon \le \frac{p}{8}$, then $|E_{g}|\ge \frac{1}{4}n^{2}p-|E_{b}| \ge\frac{1}{8}n^{2}p$.
Further, if $\varepsilon < \frac{\beta c_0 c_1 p^3}{3 \cdot 256 \cdot 64 \cdot 32}$, then since $\bar{\delta} \neq 0$ and $|\Eb| \leq \eps n^2$, we have
\[
	\sum_{ij\in E(K_{n})}\eta_{ij}
	\ge\frac{1}{3n^{2}}\cdot\frac{\beta c_{0} \muinf \overline{\delta}}{256} |E_{g}|^2
	\ge\frac{\beta c_{0} p^2}{3\cdot 256 \cdot 64}\muinf \overline{\delta} n^2
	>\frac{32}{c_{1}p}\muinf\overline{\delta}|E_{b}|.
\]
By Lemma \ref{lem:transference}, if $\varepsilon < \frac{c_1p^2}{8}$, then this implies 
\[
\sum_{ij\in E_{g}}\eta_{ij}\ge\frac{c_{1}p}{16}\sum_{ij\in E(K_{n})}\eta_{ij} > 2\muinf\overline{\delta}|E_{b}|.
\]
Therefore from (\ref{eq:gain3}), (\ref{eq:badtwist}), and Condition 3, 
if $\varepsilon \le \min\{\frac{c_0 p}{256}, \frac{c_1p^2}{8}, \frac{p}{8}, \frac{\beta c_0 c_1 p^3}{3 \cdot 256 \cdot 64 \cdot 32} \} $, then
\begin{eqnarray*}
R(T) & \ge & R(\Tnot)+\sum_{ij\in E_{g}}\eta_{ij}-\sum_{ij\in E_{b}}(|\delta_{ij}|\|t_{ij}^{(0)}\|_2+\eta_{ij})\\
 & > & R(\Tnot)+2\muinf\overline{\delta}|E_{b}|-\sum_{ij\in E_{b}}2|\delta_{ij}|\|t_{ij}^{(0)}\|_2\ge R(\Tnot).
\end{eqnarray*}
This condition on $\eps$ is satisfied under the assumption $\eps \leq \frac{\beta c_0 c_1^2 p^4}{3\cdot 256 \cdot 64 \cdot 32}$.
\end{proof}









\subsection{Properties of Gaussians in high dimensions} \label{sec-gaussians-well-distributed}

In this section, we prove that i.i.d. Gaussians satisfy properties needed to establish Conditions 2, 3, and 5 in Theorem \ref{thm:mainthm}. We begin by recording some useful facts regarding concentration of random Gaussian vectors:

\begin{lemma}
\label{lem:conc1}
Let $x,y$ be i.i.d. $\mathcal{N}(0,I_{d\times d})$, and $\epsilon \leq 1$, then
\[
\P\left ( d(1-\epsilon) \leq \|x\|_2^2 \leq d(1+\epsilon) \right) \geq 1- e^{-c \epsilon^2 d}
\]
and
\[
\P \left( |\langle x, y \rangle| \geq d\epsilon \right) \leq e^{-c\epsilon^2 d}
\]
where $c>0$ is an absolute constant.
\end{lemma}
\begin{proof}
Both statements follow from Corollary 5.17 in \cite{Vershynin}, concerning concentration of sub-exponential random variables.
\end{proof}

\begin{lemma}[\cite{Vershynin} Corollary 5.35]
\label{lem:conc2}
Let $A$ be an $n \times d$ matrix with i.i.d. $\mathcal{N}(0,1)$ entries. Then for any $t\geq 0$,
\[
\P\left( \sigma_{\max}(A) \geq \sqrt{n} + \sqrt{d} + t \right) \leq 2e^{-\frac{t^2}{2}}
\]
where $\sigma_{\max}(A)$ is the largest singular value of $A$.
\end{lemma}

%\begin{lemma}
%\label{lem: conc3}
%Let $\toi, i \in [n]$ be i.i.d. $\mathcal{N}(0,\Id)$. Then for distinct $i,j,k$,
%\[
%\P\left( \frac{99}{100}d \leq \frac{1}{2}\|t_{ij}\|_2^2 \leq \frac{101}{100}d \right) \geq 1 - e^{-cd}
%\]
%\[
%\P\left( \langle \toij,\toik \rangle^2 \geq 1/3 \right) \leq 6 e^{-cd}
%\]
%where $c>0$ is an absolute constant. 
%\end{lemma}
%\begin{proof}
%This follows from repeated application of Lemma \ref{lem:conc1}.
%\end{proof}

\begin{lemma}
\label{lem:conc3}
Let $\toi, i \in [n]$ be i.i.d. $\mathcal{N}(0,I_{d \times d})$. Then, there exists an event $E$, such that on E, we have for all $i,j,k,l \in [n], i \neq j, k\neq l$,
\[
\frac{\|\toij\|_2}{\|\tokl\|_2} \geq \frac{9}{10}
\]
and for all distinct $i,j,k \in [n]$,
\[
\langle \toijhat,\toikhat \rangle^2 \leq 1/3
\]
and $\P(E^c) \leq 3n^2 e^{-cd}$,
%\[
%\P\left( \frac{99}{100}d \leq \frac{1}{2}\|t_{ij}\|_2^2 \leq \frac{101}{100}d \right) \geq 1 - e^{-cd}
%\]
where $c>0$ is an absolute constant. 
\end{lemma}
\begin{proof}
This follows from repeated application of Lemma \ref{lem:conc1} with $\epsilon = 1/100$ and a union bound.
\end{proof}

We can now show that gaussian vectors have the well-distributed property with high probability.  Recall that $S(x,y) = \Span(x,y)$.  

\begin{lemma} \label{lem:GT}
Let $t_1,\ldots t_n \in \mathbb{R}^d$ be i.i.d. $\mathcal{N}(0,\Id)$, and let $n\geq 16$ and $d\geq 3$.  For a fixed $k \neq l$, the inequality
\[
\sum_{i \in [n], i \neq l,k} \| P_{S(t_l-t_i, t_k-t_i)^\perp}(h) \|_2 \geq \frac{1}{5}(n-2)\|P_{S(t_l-t_k)^\perp}(h)\|_2
\]
holds for all $h \in \mathbb{R}^d$ with probability of failure at most $5n e^{-cd}$, where $c>0$ is an absolute constant.
\end{lemma}

%\textcolor{blue}{Does this sound ok?}

%Note: Lemma \ref{lem:GT} immediately implies that for $T = \{t_1,\ldots, t_n \} \in \mathbb{R}^d$  with $t_i$ i.i.d. $\mathcal{N}(0,I_{d\times d})$ and further satisfying the assumptions of the lemma, we have that $T$ is well-distributed along $K_n$ with parameter $1/5$ with high probability. Using lemma \ref{lem:triangles}, this gives that $T$ is well-distributed along subgraphs $G$, with the appropriate parameter depending on the minimum number of triangles per edge in $E(G)$. 

\begin{proof}
Throughout the proof, constants named $c$ may be different from line to line, but are always bounded below by a positive absolute constant. For a fixed $(l,k)$, let $x = t_l, y = t_k$. We would like to show
\[
\sum_{i=1}^n \| P_{S(x-t_i, y-t_i)^\perp}(h) \|_2 \geq \frac{1}{5}n\|P_{S(x-y)^\perp}(h)\|_2
\]
%Where $S(v_1,v_2..v_m)$ denotes the span of vectors $v_1,\ldots v_m$ and $x,y,t_i$ satisfy certain properties to be defined. 
We note that $S(x-t_i,y-t_i) = S(x-y, x+y -2t_i)$. Thus,
\[
P_{S(x-t_i, y-t_i)^\perp}(h) = P_{S(x-y, x+y-2t_i)^\perp}(h) = P_{S(x-y, x+y-2t_i)^\perp}(P_{S(x-y)^\perp}(h))
\]
Thus, it's enough to show
\[
\sum_{i=1}^n \| P_{S(x-y, x+y-2t_i)^\perp}(h) \|_2 \geq \frac{1}{5} n\|h\|_2
\]
for $h \perp (x-y)$. 

Now, for any vectors $v,w$, we have
\[
S(v,w) = S(v, w_{v^\perp})
\]
where $w_{v^\perp} = w - \langle w,\hat{v}\rangle \hat{v}$. If $h \perp v$, we have
\begin{align*}
P_{ S(v,w)^\perp}(h) &= P_{S(v, w_{v^\perp})^\perp}(h)\\
& = h - \langle h,\hat{v}\rangle \hat{v} - \langle h, \hat{w}_{v^\perp} \rangle \hat{w}_{v^\perp}\\
&= h - \left\langle h, \frac{w}{\|w_{v^\perp}\|_2} \right\rangle \frac{w_{v^\perp}}{\|w_{v^\perp}\|_2}\\
& = h - \langle h, \hat{w} \rangle \hat{w} + \langle h, \hat{w} \rangle \left[ \hat{w} - \frac{\|w\|_2}{\| w_{v^\perp} \|_2} \frac{w_{v^\perp}}{\|w_{v^\perp}\|_2} \right]\\
& = P_{S(w)^\perp}(h) + \langle h, \hat{w} \rangle z
\end{align*}
Where $z = \hat{w} - \frac{\|w\|_2}{\| w_{v^\perp} \|_2} \frac{w_{v^\perp}}{\|w_{v^\perp}\|_2}$. Now, assuming that $|\langle \hat{v} ,\hat{w} \rangle | < 1/2$ and using that
\[
\|w_{v^\perp}\|_2^2 = \|w - \langle w, \hat{v} \rangle \hat{v} \|_2^2 = \|w\|_2^2 - 2\|w\|_2^2 \langle \hat{w}, \hat{v} \rangle + \|w\|_2^2 |\langle \hat{w},\hat{v} \rangle|^2 \geq \|w\|_2^2(1 - 2 |\langle \hat{w}, \hat{v}\rangle|)
\]
we have
\begin{align*}
\|z\|_2 &= \left\| \hat{w} - \frac{\|w\|_2}{\|w_{v^\perp}\|_2^2} \left( w - \langle w, \hat{v} \rangle \hat{v}  \right) \right\|_2\\
&= \left\| \hat{w}\left[ 1 - \frac{\|w\|_2^2}{\|w_{v^\perp}\|_2^2}\right] + \frac{\|w\|_2^2}{\|w_{v^\perp}\|_2^2}\langle \hat{w}, \hat{v} \rangle \hat{v}  \right\|_2\\
& \leq \left | 1 - \frac{\|w\|_2^2}{\|w_{v^\perp}\|_2^2} \right | + \frac{\|w\|_2^2}{\|w_{v^\perp}\|_2^2} |\langle \hat{w}, \hat{v} \rangle | = \epsilon( \langle \hat{w},\hat{v}\rangle )\\
& = \frac{\|w\|_2^2}{\|w_{v^\perp}\|_2^2}\left(1 + |\langle \hat{w}, \hat{v} \rangle | \right) - 1\\
& \leq \frac{3 |\langle \hat{w}, \hat{v} \rangle |}{1- 2|\langle \hat{w}, \hat{v} \rangle |} \triangleq \zeta( \langle\hat{w},\hat{v} \rangle )
\end{align*}
Thus, we have
\[
\| P_{S(v,w)^\perp}(h) \|_2 \geq \|P_{S(w)^\perp}(h)\|_2 -  \zeta( \langle \hat{w},\hat{v} \rangle) \|h\|_2
\]
Therefore, by taking $v = x-y$ and $w = x+y -2t_i$, to conclude the desired statement of the present Lemma, it suffices to show that
\[
\sum_{i=1}^n \| P_{S(x+y-2t_i)^\perp}(h)\|_2 \geq \gamma n\|h\|_2
\]
where $\gamma > 1/5 + \zeta \left(\frac{\langle x-y, x+y + 2t_i \rangle}{\|x-y\|_2 \| x+y -2t_i \|_2}\right)$. Note that $x-y$ and $x+y - 2t_i$ are independent, and $\frac{1}{2}(x-y) =^d \frac{1}{6}(x+y-2t_i) =^d \mathcal{N}(0,I_{d\times d})$. Applying Lemma \ref{lem:conc1} to $x-y$ and $x+y-2t_i$ with a small enough value of $\epsilon$ to ensure $\zeta\left(\frac{\langle x-y, x+y + 2t_i \rangle}{\|x-y\|_2 \| x+y -2t_i \|_2}\right) < 1/20$, we get

\[
\P\left( \zeta\left(\frac{\langle x-y, x+y + 2t_i \rangle}{\|x-y\|_2 \| x+y -2t_i \|_2}\right) > \frac{1}{20} \right) \leq 3e^{-c d}
\]
Thus, it suffices to show with high probability, that
\[
\sum_{i=1}^n \| P_{S(x+y-2t_i)^\perp}(h)\|_2 \geq 0.3 n\|h\|_2,
\]
which we proceed to establish below. 

To begin, redefine $v,w$ as $v = x+y$ and $w = -2t_i$ and consider

\begin{align*}
\sum_{i=1}^n \left \| P_{S(v+w_i)^\perp}(h)\right \|_2 &\geq \left \|  \sum_{i=1}^n P_{S(v+w_i)^\perp}(h)\right\|_2\\
&= \left \| \sum_{i=1}^n \left(h - \frac{1}{\|v+w_i\|_2^2}\langle h, v+ w_i \rangle (v+ w_i) \right)\right \|_2\\
&\geq n\|h\|_2 - \left \| \sum_{i=1}^n \frac{1}{\|v+w_i\|_2^2} (v+ w_i)(v+ w_i)^* h\right \|_2  \\
&\geq n\|h\|_2 - \left \| \sum_{i=1}^n \frac{1}{\|v+w_i\|_2^2} (v+ w_i)(v+ w_i)^* \right\|_{\text{op}} \|h\|_2\\
& \geq \|h\|_2 \left[n - \frac{1}{\min_{i}\|v+w_i\|_2^2} \left \| \sum_{i=1}^n (v+ w_i)(v+ w_i)^* \right\|_{\text{op}} \right]
\end{align*}
where in the last inequality we used 
\[
\sum_{i=1}^n \frac{1}{\|v+w_i\|_2^2} (v+ w_i)(v+ w_i)^* \preceq \frac{1}{\min_{i}\|v+w_i\|_2^2}  \sum_{i=1}^n (v+ w_i)(v+ w_i)^*
\]
Now, let $A = \sum_{i=1}^n e_i w_i^* \in \mathbb{R}^{n\times d}$. We have  
\begin{align*}
\left \| \sum_{i=1}^n (v+ w_i)(v+ w_i)^* \right\|_{\text{op}} &= \left \| \sum_{i=1}^n (vv^* + vw_i^* + w_i v^* + w_i w_i^*) \right\|_{\text{op}}\\
& \leq n\|vv^*\|_{\text{op}} + \left\| v \left( \sum_{i=1}^n w_i \right)^* +  \left( \sum_{i=1}^n w_i \right)v^* \right\|_{\text{op}} + \left \| \sum_{i=1}^n w_i w_i^* \right\|_{\text{op}}\\
& \leq n \|v\|_2^2 + 2\|v\|_2\left \| \sum_{i=1}^n w_i \right\|_2 + \left \| \sum_{i=1}^n w_i w_i^* \right\|_{\text{op}}\\
& = n \|v\|_2^2 + 2\|v\|_2\left \| \sum_{i=1}^n w_i \right\|_2 + \sigma_{\max}(A)^2\\
\end{align*}

Thus,
\[
\sum_{i=1}^n \left \| P_{S(v+w_i)^\perp}(h)\right \|_2 \geq \|h\|_2\left[n - \frac{n \|v\|_2^2 + 2\|v\|_2\left \| \sum_{i=1}^n w_i \right\|_2 + \sigma_{\max}(A)^2}{\min_{i}\|v+w_i\|_2^2} \right]
\]

Now, consider the event
\[
E = \left\{ \min_{i} \|v+w_i\|_2^2 \geq 6d\beta_1, \quad \|v\|_2^2 \leq 2d\beta_2, \quad \left \| \sum_{i=1}^n w_i \right\|_2^2 \leq 4nd\beta_3, \quad \sigma_{\max}(A)^2 \leq n\beta_4  \right\}
\]
On E, we have
\begin{align*}
\sum_{i=1}^n \left \| P_{S(v+w_i)^\perp}(h)\right \|_2 &\geq \|h\|_2 \left[ n - \frac{1}{6d\beta_1}\left( 2nd\beta_2 + 2\sqrt{2d\beta_2}2\sqrt{nd}\sqrt{\beta_3} + n\beta_4  \right) \right]\\
&= \|h\|_2 \left[n - \frac{1}{3}n \frac{\beta_2}{\beta_1} - \frac{4\sqrt{2}d\sqrt{n} \sqrt{\beta_2\beta_3}}{6d\beta_1}-\frac{\beta_4 }{6d\beta_1}n \right]\\
& = \|h\|_2 \left[n\left(1 - \frac{1}{3} \frac{\beta_2}{\beta_1} - \frac{\beta_4}{6d\beta_1}- \frac{1}{\sqrt{n}}\frac{ 4\sqrt{2} \sqrt{\beta_2\beta_3}}{6\beta_1} \right)\right]\\
\end{align*}

Now, note that $ \frac{1}{6}\|v+w_i\|_2^2 =^d \frac{1}{2}\|v\|_2^2 =^d \frac{1}{4n} \left \| \sum_{i=1}^n w_i \right\|_2^2 =^d \chi^2(d)$ and $A$ is a random $n\times d$ matrix with i.i.d. $\mathcal{N}(0,1)$ entries. 

Thus by applying Lemma \ref{lem:conc1} we have 
\[
\P\left( 6d(1-\epsilon) \leq \|v+w_i\|_2^2 \leq 6d(1+\epsilon) \right) \geq 1 - e^{-c \epsilon^2 d}
\]

\[
\P\left( 2d(1-\epsilon) \leq \|v\|_2^2 \leq 2d(1+\epsilon) \right) \geq 1 - e^{-c \epsilon^2 d}
\]

\[
\P\left( 4nd(1-\epsilon) \leq \left \| \sum_{i=1}^n w_i \right\|_2^2 \leq 4nd(1+\epsilon) \right) \geq 1 - e^{-c \epsilon^2 d}
\]

where $c>0$ is a universal constant. Also by taking $t = \sqrt{2d}$ in Lemma \ref{lem:conc2} we get
\[
\P \left( \sigma_{\text{max}}(A) \geq \sqrt{n} +2\sqrt{d} \right) \leq 2e^{-d}
\]

Now, let $\beta_1 = 1-\frac{1}{100}, \quad \beta_2 = \beta_3 = 1 + \frac{1}{100}, \quad \beta_4 = \frac{d}{2}$, which gives
\[
\frac{1}{3} \frac{\beta_2}{\beta_1} \leq 1/3 +1/99, \quad \frac{1}{\sqrt{n}}\frac{ 4\sqrt{2} \sqrt{\beta_2\beta_3}}{6\beta_1} < \frac{1}{\sqrt{n}}, \quad \frac{\beta_4}{5d} = 1/10
\]

We have
\[
\P  \left(\sigma_{\max} (A) \geq \sqrt{n \beta_4}\right) \leq \P\left( \sigma_{\max} (A) \geq \sqrt{n} + 2\sqrt{d} \right) \leq 2e^{-d}
\]
whenever $\sqrt{n} + 2\sqrt{d} \leq \sqrt{n}\sqrt{d/2}$, which holds whenever
\[
n \geq \left( \frac{2\sqrt{d}}{\sqrt{d/2} -1} \right)^2
\]
which holds for $n \geq 16$ when $d\geq 3$. Since for $n\geq 16$, $\frac{1}{\sqrt{n}} \leq 1/4$, we have on E
\[
\sum_{i=1}^n \left \| P_{S(v+w_i)^\perp}(h)\right \|_2 \geq 0.3 n \|h\|_2
\]
Thus,
\[
\P\left ( \sum_{i=1}^n \left \| P_{S(v+w_i)^\perp}(h)\right \|_2 < 0.3 n \|h\|_2 \right) \leq \P(E^c) \leq (n+3)e^{-cd}
\]
where $c>0$ is an absolute constant. 

Combining all of the above, we get
\[
\sum_{i=1}^n \| P_{S(x-t_i, y-t_i)^\perp}(h) \|_2 \geq \frac{1}{5}n\|P_{S(x-y)^\perp}(h)\|_2
\]

with probability of failure at most $5ne^{-cd}$. 

%Finally, recalling that $x = t_l, y= t_k$ and applying a union bound over $(l,k) \in [n] \times [n], l \neq k$, we have
%\[
%\sum_{i \in [n], i \neq l,k} \| P_{S(t_l-t_i, t_k-t_i)^\perp}(h) \|_2 \geq \frac{1}{5}(n-2)\|P_{S(t_l-t_k)^\perp}(h)\|_2
%\]
%for any $(l,k) \in [n]\times [n], l \neq k$, with probability at least $1- 5n^3 e^{-cd}$.

 \end{proof}


\begin{lemma} \label{lem:ptyp-well}
Let $G([n],E)$ be $p$-typical, and $t_1,\ldots t_n \sim \mathcal{N}(0,I_{d\times d})$ be i.i.d. Then $T = \{t_i\}_{i \in [n]}$ is $\frac{1}{5}$-well distributed along $G$ with probability at least $1-10n^3 e^{-cd}$, where $c>0$ is an absolute constant. 
\end{lemma}
\begin{proof}
For each $ij \in E$, let $S_{ij} = \{ k \in [n]; ik, jk \in E(G) \}$ and note that $|S_{ij}| \leq 2np^2$. Now apply Lemma \ref{lem:GT} to the set of vectors $ \{t_i, t_j\} \bigcup \{ t_k \}_{k \in \mathcal{I}_{ij}}$, with the distinguished vectors being $\{t_i, t_j\}$, which gives the desired property for the pair $(i,j)$ with probability of failure at most $5(|S_{ij}|)e^{-cd} \leq 5 (2np^2) e^{-cd}$, where $c>0$ is an absolute constant. Taking the union bound over pairs of distinct integers $i,j \in [n]$, we get the desired property simultaneously for all pairs with probability at least $1- n^2 \cdot 5(2np^2) e^{-cd} = 1 - 10n^3 p^2 e^{-cd} \geq 1- 10n^3 e^{-cd}$.
\end{proof}

%%
%%
%%
%%	SECTION : Random graphs
%%
%%
%%
\subsection{Random graphs are $p$-typical with high probability} \label{sec-random-graph}



\begin{lemma} \label{lem:ptyp}
There exists an absolute constant $c > 0$ such that for all positive real numbers $p \le 1$, $G(n,p)$ is $p$-typical with probability at least $1-n^2e^{-cnp^2}$ if $np \ge 4 \log n$.
\end{lemma}
\begin{proof}
A graph is not connected only if there exists a partition $V_1 \cup V_2$ of its vertex
set for which there are no edges between $V_1$ and $V_2$. Without loss of generality,
we may assume that $|V_1| \le \lfloor \frac{n}{2} \rfloor$. Since the number
of ways to choose a set of size $k$ from a set of size $n$ is ${n \choose k}$, 
the probability that $G(n,p)$ is not connected is at most 
\[
	\sum_{k=1}^{\lfloor n/2 \rfloor} {n \choose k} (1-p)^{k(n-k)}
	\le \sum_{k=1}^{\lfloor n/2 \rfloor} \left(\frac{en}{k}\right)^{k} e^{-pk(n-k)}
	< \sum_{k=1}^{\lfloor n/2 \rfloor} \left(n e^{1-p(n-k)} \right)^{k}.
\]
Since $k \le \lfloor \frac{n}{2} \rfloor$, we have
$n e^{1-p(n-k)} \le n e^{1-pn/2} < 1$ (since $np \ge 4 \log n$). 
Therefore the summand on the right-hand-side is at most $(ne^{1-pn/2})^{k}$,
which is maximized at $k=1$. 
This shows that the probability that $G(n,p)$ is not connected is at most $n^2 e^{1-pn/2}$.

In $G(n,p)$, for a fixed vertex $v$, the expected value of $\deg(v)$is
$(n-1)p$, and for a pair of vertices $v,w$, the expected value of
the codegree of $v$ and $w$ is $(n-2)p^{2}$. Therefore the lemma
follows from Chernoff's inequality --- see Fact 4 from  \cite{AngluinValiant} --- and a union bound. 
\end{proof}





%%
%%
%%
%%	SECTION : main
%%
%%
%%

%\begin{theorem} \label{thm:mainthm}
%For all $p,\beta,c_{0}$ and $c_{1}$, there exists
%$\varepsilon$ such that under the assumptions stated above, we have
%$R(T)>R(\Tnot)$ for all $T\neq \Tnot$ such that $L(T)=1$.
%\end{theorem}


\subsection{Proof of Theorem \ref{thm-complete}} \label{sec-proof-random-theorem}

We can now prove the high dimensional recovery theorem, which we state here again for convenience:

\addtocounter{theorem}{-3}

\begin{theorem} 
Let $G([n],E)$ be drawn from $G(n,p)$ for some $p = \Omega(n^{-1/4})$. Take $ \tnot_1, \ldots \tnot_n \sim \mathcal{N}(0, I_{d \times d})$ to be i.i.d., independent from $G$. There exists an absolute constant $c>0$ and a $\gamma = \Omega(p^4)$ not depending on $n$, such that if $\max(\frac{2^6}{c^6}, \frac{4^3}{c^3} \log^3 n) \leq  n \leq e^{\frac{1}{6}c d}$ and $d = \Omega(1)$, then there exists an event with probability at least  $1- e^{-n^{1/6}} - 13 e^{-\frac{1}{2}c d}$, on which the following holds:\\[1em]
For arbitrary  subgraphs $\Eb$ satisfying $\max_i \deg_b(i) \leq \gamma n$ and arbitrary pairwise direction corruptions  $\vij \in \mathbb{S}^{d-1}$ for $ij \in \Eb$,  the convex program \eqref{shapefit} has a unique minimizer equal to $\left \{\alpha \Bigl(\toi - \tnotbar \Bigr)\right\}_{i \in [n]}$ for some positive $\alpha$ and for $\tnotbar = \frac{1}{n}\sum_{i\in [n]} \toi$. 
\end{theorem}
\begin{proof}
It is enough to verify that $G$, $T$ and $E_b$ in the assumption of the present theorem satisfy the deterministic conditions 1--6 in Theorem 2, with appropriate constants $p, \beta, c_0, \epsilon, c_1$, and with the purported probability. By Lemma \ref{lem:ptyp}, Lemma \ref{lem:conc3}, and Lemma \ref{lem:ptyp-well}, we have that Condition 1 holds with value $p$, Condition 2 holds with $\beta = \sqrt{\frac{2}{3}}$, Condition 3 holds with $c_0 = \frac{9}{10}$, and Condition 5 holds with $c_1 = \frac{1}{5}$, with probability at least
\[
1-n^2 e^{-cnp^2} - 3n^2 e^{-cd} - 10n^3 e^{-cd}
\]
where $c>0$ is an absolute constant.

Thus, taking any $E_b$, which satisfies Condition 4 with $\gamma = \frac{p^4}{10^7} \leq \frac{\beta c_0 c_1^2 p^4}{256\cdot 32\cdot 64 \cdot 3}$ , we get that recovery via ShapeFit is guaranteed. Note that the condition $\max \deg_b(i) \leq \gamma n$ is nontrivial when $p = \Omega(n^{-1/4})$.  Using the requirements on $n$ and $p$, we have $n^2 e^{-cnp^2} \leq n^2 e^{-cn^{1/3}} \leq e^{-\frac{1}{6}n}$ and $13n^3 e^{-cd} \leq 13(e^{\frac{1}{6}cd})^3 e^{-cd} \leq 13 e^{-\frac{1}{2}cd}$. Thus, the probability of exact recovery via ShapeFit, uniformly in $E_b$ and $v_{ij}$ satisfying the assumptions of the theorem, is at least
\[
1- e^{-n^{1/6}} - 13 e^{-\frac{1}{2}c d}. \qedhere
\]

%probability at least, that (1) holds with probability at least... Then the conclusion of the theorem guarantees that taking $\gamma < \frac{\beta c_0 c_1^2 p^4}{10^6}$, guarantees recovery via ShapeFit. 


\end{proof}





































































































\section{Proof of three-dimensional recovery} \label{sec-proofs-three-d}

The proof of recovery in three dimensions parallels the proof in high dimensions, but it is more technical because it can not capitalize on the concentration of measure phenomenon in high dimensions.
Specifically, the additional technicality in three dimensions comes from the fact that
for large $n$, there exist pairs of locations $\toi$, $\toj$
that are close to each other, i.e., $\|\toij\|_2$ is small.
For such pairs of vectors, with high probability, for all $k \neq i,j$
the value of $1 - \langle \toikhat, \tojkhat \rangle^2$ will be small.
% (in fact proportional to $\|t_{ij}^{(0)}\|_2^2$).
This fact introduces the following two main obstacles in carrying out the same analysis:
\begin{enumerate}
  \setlength{\itemsep}{1pt} \setlength{\parskip}{0pt}
  \setlength{\parsep}{0pt}
  \item There is no uniform lower bound on  $1 - \langle \toikhat, \tojkhat \rangle^2$.  Hence Condition 2 in Theorem~\ref{thm:mainthm} fails.
  \item There is no uniform lower bound on $\|t_{ij}^{(0)}\|_2$. Hence Condition 3 in Theorem~\ref{thm:mainthm} fails. 
\end{enumerate}
These are indeed obstacles since
the gains in rotational motions coming from Lemmas~\ref{lem:rigidity1} and \ref{lem:rigidity2}
are proportional to $\sqrt{1 - \langle \toikhat, \tojkhat \rangle^2}$ and $\| \toij \|_2$.
We avoid these difficulties and prove the three-dimensional analogue of Theorem~\ref{thm:mainthm} by weakening Conditions 2 and 3. 
Roughly speaking, in $\mathbb{R}^3$, Condition 2 holds for most triples $i,j,k \in [n]$ (instead of all triples)
and Condition 3 gets replaced by a one-sided version where we only have a uniform  upper bound on the lengths $\| \toij \|_2$.

%One can easily check that for large $n$, there will be pairs of locations $t_{i}^{(0)}$, $t_{j}^{(0)}$
%that are close to each other, forcing $\|t_{ij}^{(0)}\|$ to be small (thus the second obstacle). 
%For these pairs, the gain on 
%rotational motions obtained by applying \red{Lemmas~\ref{lem:rigidity1} and \ref{lem:rigidity2}} would be rather %small since the gain is
%proportional to $\|t_{ij}^{(0)}\|_2$. 
%Furthermore, for such pairs of vectors, with high probability, for all $k \neq i,j$
%we will have $1 - \langle t_{ik}, t_{jk} \rangle^2$ small (in fact proportional to $\|t_{ij}^{(0)}\|_2^2$).
%Fortunately, \red{Lemmas~\ref{lem:rigidity1} and \ref{lem:rigidity2}} can still be applied to these pairs since the lemma
%only requires a lower bound on $1 - \langle t_{ij}, t_{ik} \rangle^2$ and 
%$1 - \langle t_{ij}, t_{jk} \rangle^2$. 
%Thus we prove \red{the three-dimensional analogue of Theorem~\ref{thm:mainthm}} by showing that Conditions 2 and 3 are valid for most of the vectors $t_{ij}^{(0)}$. 

Unlike in the high-dimensional case where we allowed a constant fraction of edges incident to each
vertex to be corrupted, the three-dimensional case requires the fraction of corrupted edges incident 
to each vertex to be at most $O(\frac{1}{\log^3 n})$.
This additional poly-logarithmic factor is due to the fact that our well-distributedness proof in three dimensions hinges on the maximum $\ell_2$ norm of locations, which is $\Omega(\sqrt{\log n})$ with high probability.
It can be removed for a distribution of locations that has a uniform constant upper bound on $\|\toi\|_2$.


\subsection{Deterministic recovery theorem in three dimensions} \label{sec-deterministic-theorem-three-d}

%\blue{(Note: omitted definition of general positions)}
%We say that a set of vectors in $\mathbb{R}^3$ is in {\em general position} if
%no \red{plane} contains  four of the vectors in the set.  \red{(Do we really need general position? or is it enough to just have $\toij \neq 0$?)} \blue{[Case 0 uses the fact that $\beta\neq 0$ for all triples. Since $\beta=0$ if and only if $t_i, t_j, t_k$ are co-linear, we need them to be in general position.]}

We now state deterministic conditions on the graph $G$, the corrupted observations $\Eb$, and the locations $\Tnot$ that guarantee recovery.
Recall the definition $\mu=\frac{1}{|E(G)|}\sum_{ij\in E(G)}\|t_{ij}^{(0)}\|_2$. 

\setcounter{theorem}{3}
\begin{theorem} \label{thm-deterministic-three-d}
Suppose $\Tnot, \Eb, G$ satisfy the conditions
\begin{enumerate}
  \setlength{\itemsep}{1pt} \setlength{\parskip}{0pt}
  \setlength{\parsep}{0pt}
	\item \label{enu:typical}The underlying graph $G$ is $p$-typical,
	\item \label{enu:angle}For all distinct $i,j\in[n]$, for all but at most $\eps_{1}n$ indices $k\in[n]$ satisfying $k \neq i,j$, we have $1-\langle\hat{t}_{ij},\hat{t}_{ik}\rangle^{2}\ge\beta^{2}$ and $1-\langle\hat{t}_{ij},\hat{t}_{jk}\rangle^{2}\ge\beta^{2}$,
	\item \label{enu:length}For all distinct $i,j\in[n]$, we have $\|t_{ij}^{(0)}\|_2\le c_{0}\mu$, 	\item \label{enu:bad}Each vertex has at most $\eps_0 n$ edges in $\Eb$ incident to it,
	\item \label{enu:welldistributed}The set $\{t_{i}^{(0)}\}_{i \in [n]}$ is $c_{1}$-well-distributed
along $G$,
	\item \label{enu:general}No three vectors $\toi, \toj, \tok$ are collinear for distinct $i,j,k$.
	% (Isn't this enough?)} The set of vectors $\{\toi\}$ is in general position \blue{Added the definition before the statement of the theorem. We used it in Case 0. I thought we implicitly used it other places to guarantee $t_i \neq t_j$ and so on, but am not sure if we really did. If the only place it is used is in Case 0, then this condition can be removed. },
\end{enumerate}
for constants $0< p, \beta, \eps_0, \eps_1, c_{1} \leq 1 \le c_0$.  If $\eps_0 \leq \frac{\beta c_1^2 p^4}{32 \cdot 3 \cdot 64 \cdot 1024 c_0^2}$ and $\eps_1 \le \frac{p}{192c_0}$, then $L(\Tnot)\neq 0$ and $\Tnot / L(\Tnot)$ is the unique optimizer of ShapeFit.
\end{theorem}

Note that all six 
conditions are invariant under translation and non-zero scalings of $\Tnot$ (Condition~\ref{enu:length}
is invariant since both $t_{ij}^{(0)}$ and $\mu$ scale together and are invariant under translation).
Before we prove the theorem, we establish that $L(\Tnot)\neq 0$ when $\eps_0$ is small.  This non-equality guarantees that some scaling of $\Tnot$ is feasible whenever, roughly speaking, $|\Eb| < |\Eg|$.


\begin{lemma} \label{lem:T0-feasible-three-d} 
If $\eps_0 < \frac{p}{8c_0}$, then $L(\Tnot) \neq 0$.
\end{lemma}
\begin{proof}
Since $v_{ij} = \toijhat$ for all $ij \in E_g$, we have
\[
	L(\Tnot) = \sum_{ij \in E(G)} \langle \toij, v_{ij} \rangle 
	\geq \sum_{ij \in E_g} \| \toij\|_2 - \sum_{ij \in E_b} \| \toij \|_2
	= \sum_{ij \in E(G)} \| \toij\|_2 - 2\sum_{ij \in E_b} \| \toij \|_2.
\]
By Condition \ref{enu:length},  $\sum_{ij \in E_b} \| \toij \|_2 \le c_0 \mu |E_b| \le c_0\mu \cdot \eps_0n^2 < \frac{1}{8}n^2p \mu$. Since Condition \ref{enu:typical} implies $|E(G)| \ge \frac{1}{4}n^2p$, we have
$\sum_{ij \in E(G)} \| \toij\|_2 \ge \frac{1}{4}n^2p\mu$. Therefore it follows that $L(\Tnot) > 0$. 
\end{proof}






%%
%%
%%
%%	SECTION : main
%%
%%
%%


\subsection{Proof of Theorem \ref{thm-deterministic-three-d}} \label{sec-proof-deterministic-3d}

Lemmas \ref{lem:rigidity1} and \ref{lem:rigidity2} will be repeatedly
used throughout the proof. Note that these lemmas can be used only
if the given set of vectors satisfies a certain condition on the angles between
them. For each distinct $ij\in E(K_{n})$, define $B(ij)$ as the
set of edges $k\ell\in E(K_{n})$ such that $\sqrt{1-\langle\hat{t}_{ac}^{(0)},\hat{t}_{bc}^{(0)}\rangle^{2}}<\beta$
holds for some distinct $a,b,c\in\{i,j,k,\ell\}$ satisfying
$(a,b)\neq(i,j)$. Note that Lemmas~\ref{lem:rigidity1} and \ref{lem:rigidity2} can be applied to the
set of indices $\{i,j,k,\ell\}$ (having size either 3 or 4) for all $k\ell \notin B(ij)$.
The following lemma shows that $B(ij)$ is small for each $ij$.


\begin{lemma}
\label{lem:badK4}For each $ij\in E(K_{n})$, we have $|B(ij)|\le6\varepsilon_{1}n^{2}$.
\end{lemma}
\begin{proof}
For each $ab\in E(K_{n})$, define $B_{3}(ab)$ as the set of indices
$c\in[n]$ distinct from $a,b$ for which $\sqrt{1-\langle \hat{t}_{ab}^{(0)}, \hat{t}_{ac}^{(0)}\rangle^{2}}<\beta$
or $\sqrt{1-\langle\hat{t}_{ab}^{(0)},\hat{t}_{bc}^{(0)}\rangle^{2}}<\beta$
holds. Condition \ref{enu:angle} implies $|B_{3}(ab)|\le\varepsilon_{1}n$
for all $ab \in E(K_n)$. One can check that $k\ell\in B(ij)$ if and only if
one of the following events hold: $k\in B_{3}(ij)$, $\ell\in B_{3}(ij),$
$k\in B_{3}(i\ell)\cup B_{3}(j\ell)$, $\ell\in B_{3}(ik)\cup B_{3}(jk)$.
Therefore 
\begin{eqnarray*}
|B(ij)| & \le & 2|B_{3}(ij)|\cdot n+\sum_{\ell\neq i,j}\Big(|B_{3}(i\ell)|+|B_{3}(j\ell)|\Big)+\sum_{k\neq i,j}\Big(|B_{3}(ik)|+|B_{3}(jk)|\Big)\\
 & \le & 2\varepsilon_{1}n^{2}+n\cdot2\varepsilon_{1}n+n\cdot2\varepsilon_{1}n=6\varepsilon_{1}n^{2}.\qedhere
\end{eqnarray*}
\end{proof}


We now prove the deterministic recovery theorem in three dimensions.

\begin{proof}[Proof of Theorem \ref{thm-deterministic-three-d}]

By Lemma \ref{lem:T0-feasible-three-d} and the fact that Conditions 1--6 are invariant under global translation and nonzero scaling, we can take $\Tnotbar=0$ and $L(\Tnot) = 1$ without loss of generality.  The variable $\mu$ is to be understood accordingly.  

We will directly prove that $R(T) > R(\Tnot)$ for all $T\neq \Tnot$ such that $L(T)=1$ and $\Tbar = 0$. Consider an arbitrary feasible $T$ and recall the notation $\tij = (1 + \deltaij) \toij + \etaij \sij$ where $\sij$ is a unit vector orthogonal to $\toij$ and $\etaij = \| P_{\toijperp} \tij \|_2$. A useful lower bound for the objective $R(T)$ is given by 
\begin{eqnarray}
	R(T)
	=\sum_{ij \in E(G)}\|\proj_{v_{ij}^{\perp}}t_{ij}\| _2
	& = & \sum_{ij\in E_{g}}\eta_{ij}+\sum_{ij\in E_{b}}\|\proj_{v_{ij}^{\perp}}t_{ij}\|_2\nonumber \\
	& \ge & \sum_{ij\in E_{g}}\eta_{ij}+\sum_{ij\in E_{b}}\left(\|\proj_{v_{ij}^{\perp}}t_{ij}^{(0)}\|_2-|\delta_{ij}|\|t_{ij}^{(0)}\|_2-\eta_{ij}\right)\nonumber \\
	& = & R(\Tnot)+\sum_{ij\in E_{g}}\eta_{ij}-\sum_{ij\in E_{b}}(|\delta_{ij}|\|t_{ij}^{(0)}\|_2+\eta_{ij}).\label{eq:gain3-3d}
\end{eqnarray}

Suppose that $\sum_{ij\in E_{b}}|\delta_{ij}|\|t_{ij}^{(0)}\|_2<\sum_{ij\in E_{b}}\eta_{ij}$.
Since $\varepsilon_0 \le\frac{c_{1}p^{2}}{16}$, Lemma \ref{lem:badgood}
implies $\sum_{ij\in E_{b}}\eta_{ij}\le\frac{1}{2}\sum_{ij\in E_{g}}\eta_{ij}$.
Therefore by (\ref{eq:gain3-3d}), we have
\begin{eqnarray*}
R(T) & \ge & R(T_{0})+\sum_{ij\in E_{g}}\eta_{ij}-\sum_{ij\in E_{b}}(|\delta_{ij}|\|t_{ij}^{(0)}\|_2+\eta_{ij})\\
 & > & R(T_{0})+\sum_{ij\in E_{g}}\eta_{ij}-\sum_{ij\in E_{b}}2\eta_{ij}\ge R(T_{0}).
\end{eqnarray*}
Hence we may assume 
\begin{equation}
\sum_{ij\in E_{b}}|\delta_{ij}|\|t_{ij}^{(0)}\|_2\ge\sum_{ij\in E_{b}}\eta_{ij}.\label{eq:badtwist-3d}
\end{equation}
In other words, the total parallel motion is larger than the total rotational
motions on the bad edges. The key idea of the proof is to show that
parallel motions on bad edges induce a large amount of rotational motions
on good edges. 

In the case $|\Eb| \neq 0$,  define $\overline{\delta}:=\frac{\sum_{ij\in E_{b}}|\delta_{ij}|\|t_{ij}^{(0)}\|_2}{\sum_{ij\in E_{b}}\|t_{ij}^{(0)}\|_2}$
as the average `relative parallel motion' on the bad edges. 
For distinct $ij,k\ell\in E(K_{n})$, if $\{i,j\}\cap\{k,\ell\}=\emptyset$,
then define $\eta(ij,k\ell)=\eta_{ij}+\eta_{ik}+\eta_{i\ell}+\eta_{jk}+\eta_{j\ell}+\eta_{k\ell}$,
and if $\{i,j\}\cap\{k,\ell\}\neq\emptyset$ (without loss of generality,
assume $\ell=i$), then define $\eta(ij,k\ell)=\eta_{ij}+\eta_{ik}+\eta_{jk}$. 

\medskip{}

\noindent \textbf{Case 0}. $\bar{\delta} = 0$ or $|\Eb| = 0$.
\medskip

Note that $\bar{\delta}=0$ implies $\delta_{ij} =0$ for all $ij \in E_b$, which by \eqref{eq:badtwist} implies
$\eta_{ij} =0$ for all $ij \in E_b$.  Therefore by \eqref{eq:gain3}, we have
\[
	R(T) \geq R(\Tnot) + \sum_{ij \in E_g} \eta_{ij}.
\]
If $\sum_{ij \in E_g} \eta_{ij} > 0$, then we have $R(T) > R(\Tnot)$. Thus 
we may assume that $\eta_{ij} = 0$ for all $ij \in E_g$. In this case, we will show that $T = \Tnot$.

By Lemma~\ref{lem:transference}, if $\varepsilon_0 \le \frac{c_1p^2}{8}$, then
$\eta_{ij}=0$ for all $ij \in E(G)$ implies that
$\eta_{ij} = 0$ for all $ij \in E(K_n)$. 
For $ij \in E_b$, since $\delta_{ij} = \eta_{ij} = 0$, it follows that $\ell_{ij} = \ell_{ij}^{(0)}$.
Since $\delta_{ij} = \ell_{ij} - \ell_{ij}^{(0)}$ for $ij \in E_g$, we have
\[
	0 = \sum_{ij \in E(G)} (\ell_{ij} - \ell_{ij}^{(0)})
	= \sum_{ij \in E_b} (\ell_{ij} - \ell_{ij}^{(0)}) + \sum_{ij \in E_g} (\ell_{ij} - \ell_{ij}^{(0)})
	= \sum_{ij \in E_g} (\ell_{ij} - \ell_{ij}^{(0)})
	= \sum_{ij \in E_g} \delta_{ij}\|\toij\|_2,
\] 
where the first equality is because $L(T) = L(\Tnot) = 1$.
By Condition 0, we have $\|t_{ij}^{(0)}\| \neq 0$ for all $ij \in E_g$.
%Since the vectors $\{t_i\}_{i \in [n]}$ are in general position, we have $\|t_{ij}^{(0)}\| \neq 0$ for all $ij \in E_g$. 
Hence if $\delta_{ij} \neq 0$ for some $ij \in E_g$, then there exists $ab, cd \in E_g$ 
such that $\delta_{ab} > 0$ and $\delta_{cd} < 0$. 
By Lemma~\ref{lem:rigidity1} or \ref{lem:rigidity2} and Condition 6, 
this forces $\eta(ab,cd) > 0$, contradicting the fact that
$\eta_{ij} = 0$ for all $ij \in E(K_n)$. Therefore $\delta_{ij} = 0$ for all $ij \in E_g$, and hence
$\delta_{ij} = 0$ for all $ij \in E(G)$. 

Define $t_i = t_i^{(0)} + h_i$ for each $i \in [n]$. 
Because $\eta_{ij} = \delta_{ij} = 0$ for all $ij \in E(G)$, we have $h_i = h_j$ 
for all $ij \in E(G)$. Since $G$ is connected (by Condition~\ref{enu:typical}), this implies $h_i = h_j$ for all 
$i \in [n]$. Then by the constraint $\sum_{i \in [n]} t_i = \sum_{i \in [n]} t_i^{(0)} = 0$, 
we get $h_i = 0$ for all $i \in [n]$. Therefore $T = \Tnot$.
This proves Case 0. \\

\medskip

We may now assume that $\overline{\delta}\neq0$.
Since $\ell_{ij}-\ell_{ij}^{(0)}=\delta_{ij}\|t_{ij}^{(0)}\|_2$ for
$ij\in E_{g}$, we have 
\begin{eqnarray*}
0=\sum_{ij \in E(G)}(\ell_{ij}-\ell_{ij}^{(0)}) & = & \sum_{ij\in E_{b}}(\ell_{ij}-\ell_{ij}^{(0)})+\sum_{ij\in E_{g}}\delta_{ij}\|t_{ij}^{(0)}\|_2.
\end{eqnarray*}
Therefore 
\begin{eqnarray}
\left|\sum_{ij\in E_{g}}\delta_{ij}\|t_{ij}^{(0)}\|_2\right| & \le & \left|\sum_{ij\in E_{b}}(\ell_{ij}-\ell_{ij}^{(0)})\right|\le\sum_{ij\in E_{b}}(|\delta_{ij}|\|t_{ij}^{(0)}\|_2+\eta_{ij})\nonumber \\
 & \le & 2\sum_{ij\in E_{b}}|\delta_{ij}|\|t_{ij}^{(0)}\|_2.\label{eq:good_parallel}
\end{eqnarray}
where the final inequality follows from (\ref{eq:badtwist}). 

Define $E_{g}'=\{ij\in E_{g}\,:\,\|t_{ij}^{(0)}\|_2\ge\frac{1}{2}\mu\}$
as the set of `long' good edges. Since $\sum_{ij\in E_{g}\setminus E_{g}'}\|t_{ij}^{(0)}\|_2 < \frac{1}{2}\mu|E_{g}|$,
we have 
\begin{eqnarray*}
	\sum_{ij\in E_{g}'}\|t_{ij}^{(0)}\|_2 
	& = & \sum_{ij\in E(G)}\|t_{ij}^{(0)}\|_2-\sum_{ij\in E_{b}}\|t_{ij}^{(0)}\|_2-\sum_{ij\in E_{g}\setminus E_{g}'}\|t_{ij}^{(0)}\|_2\\
	& > & \mu|E(G)|-c_{0}\mu\cdot|E_{b}|-\frac{1}{2}\mu|E_{g}| \geq \mu\cdot\frac{1}{16}n^{2}p.
\end{eqnarray*}
where the last inequality uses $|E(G)| \geq \frac{n^2 p}{4}$, $|\Eb| < \eps_0 n^2$, $|\Eg| \leq |E(G)|$, $\eps_0 < \frac{p}{16 c_0}$.
By Condition~\ref{enu:length}, we have $\|t_{ij}^{(0)}\|_2\le c_{0}\mu$
for all $ij$, and thus it follows that 
\begin{equation}\label{eq:longgoodedges}
|E_{g}'|\ge\frac{1}{16c_{0}}n^{2}p.
\end{equation}

\medskip

\noindent\textbf{Case 1}. $\overline{\delta} \neq 0$ and  $\sum_{ij\in E_{g}'}|\delta_{ij}|<\frac{1}{8}\overline{\delta}|E_{g}'|$ and $|\Eb| \neq 0$.

In this case, we will exploit the fact that there is a difference between average
relative parallel motions on long good edges and that on bad edges, to show that there is
a large amount of rotational motion on the $K_{4}$s of the form $\{i,j,k,\ell\}$
where $ij\in E_{b}$ and $k\ell\in E_{g}'$. Define $L_{b}=\{ij\in E_{b}:|\delta_{ij}|\ge\frac{1}{2}\overline{\delta}\}.$
Note that $\sum_{ij\in E_{b}\setminus L_{b}}|\delta_{ij}|\|t_{ij}^{(0)}\|_2<\frac{1}{2}\overline{\delta}\sum_{ij\in E_{b}}\|t_{ij}^{(0)}\|_2=\frac{1}{2}\sum_{ij\in E_{b}}|\delta_{ij}|\|t_{ij}^{(0)}\|_2$.
Therefore 
\begin{eqnarray}
\sum_{ij\in L_{b}}|\delta_{ij}|\|t_{ij}^{(0)}\|_2 & = & \sum_{ij\in E_{b}}|\delta_{ij}|\|t_{ij}^{(0)}\|_2-\sum_{ij\in E_{b}\setminus L_{b}}|\delta_{ij}|\|t_{ij}^{(0)}\|_2>\frac{1}{2}\sum_{ij\in E_{b}}|\delta_{ij}|\|t_{ij}^{(0)}\|_2.\label{eq:large_delta_3d}
\end{eqnarray}
Define $F_{g}=\{ij\in E_{g}:|\delta_{ij}|<\frac{1}{4}\overline{\delta}\}$.
Then by the condition of Case 1,
\[
\frac{1}{8}\overline{\delta}|E_{g}'|>\sum_{ij\in E_{g}'}|\delta_{ij}|\ge\sum_{ij\in E_{g}'\setminus F_{g}}|\delta_{ij}|\ge\frac{1}{4}\overline{\delta}|E_{g}'\setminus F_{g}|,
\]
and therefore $|E_{g}'\setminus F_{g}|<\frac{1}{2}|E_{g}'|$, or equivalently,
$|F_{g}|>\frac{1}{2}|E_{g}'|\ge\frac{1}{32c_{0}}n^{2}p$ (where the
second inequality comes from (\ref{eq:longgoodedges})).

For each $ij\in L_{b}$ and $k\ell\in F_{g}\setminus B(ij)$, by Lemmas
\ref{lem:rigidity1} and \ref{lem:rigidity2}, we have $\eta(ij,k\ell)\ge\frac{\beta}{4}|\delta_{k\ell}-\delta_{ij}|\|t_{ij}^{(0)}\|_2\ge\frac{\beta}{4}\cdot\frac{1}{2}|\delta_{ij}|\|t_{ij}^{(0)}\|_2$.
Therefore, 
\begin{eqnarray*}
\sum_{ij\in E_{b}}\sum_{k\ell\in E_{g}}\eta(ij,k\ell) & \ge & \sum_{ij\in L_{b}}\sum_{k\ell\in F_{g}\setminus B(ij)}\frac{\beta}{8}|\delta_{ij}|\|t_{ij}^{(0)}\|_2=\sum_{ij\in L_{b}}|F_{g}\setminus B(ij)|\cdot\frac{\beta}{8}|\delta_{ij}|\|t_{ij}^{(0)}\|_2.
\end{eqnarray*}
By Lemma \ref{lem:badK4}, we know that $|B(ij)|<6\varepsilon_{1}n^{2}$
holds for all $ij\in E(K_{n})$. For $\varepsilon_1 \le \frac{p}{192c_0}$, we have
\[
	|F_{g}\setminus B(ij)| 
	> \frac{1}{32c_{0}}n^{2}p-6\varepsilon_{1}n^{2}
	\ge\frac{1}{64c_{0}}n^{2}p.
\]
Therefore
\[
	\sum_{ij\in E_{b}}\sum_{k\ell\in E_{g}}\eta(ij,k\ell)
	> \frac{\beta}{8}\cdot\frac{1}{64c_{0}}n^{2}p\cdot\sum_{ij\in L_{b}}|\delta_{ij}|\|t_{ij}^{(0)}\|_2\ge\frac{\beta}{1024c_{0}}n^{2}p\sum_{ij\in E_{b}}|\delta_{ij}|\|t_{ij}^{(0)}\|_2,
\]
where the second inequality comes from (\ref{eq:large_delta_3d}).

For each $ij\in E(K_{n})$, we would like to count how many times
each $\eta_{ij}$ appear on the left hand side. If $ij\in E_{b}$,
then there are at most ${n \choose 2}$ $K_{4}$s and $n$ $K_{3}$s
containing $ij$; hence $\eta_{ij}$ may appear at most $6{n \choose 2}+3n=3n^{2}$
times. If $ij\notin E_{b}$, then $\eta_{ij}$ appears when there
is a $K_{4}$ or a $K_{3}$ containing $ij$ and some bad edge. By
Condition 4, there are at most $2\varepsilon_{0}n$ such bad $K_{3}$s.
If the bad edge in $K_{4}$ is incident to $ij$, then there are at
most $2\varepsilon_{0}n\cdot(n-3)$ such $K_{4}$s, and if the bad
edge is not incident to $ij$, then there are at most $|E_{b}|\le\varepsilon_{0}n^{2}$
such $K_{4}$s. Thus $\eta_{ij}$ may appear at most $3\cdot2\varepsilon_{0}n+6\cdot(2\varepsilon_{0}n(n-3)+\varepsilon_{0}n^{2})\le18\varepsilon_{0}n^{2}$
times. Therefore
\begin{eqnarray*}
\sum_{ij\in E_{b}}\sum_{k\ell\in E_{g}}\eta(ij,k\ell) & \le & \sum_{ij\in E_{b}}3n^{2}\cdot\eta_{ij}+\sum_{ij\in E(K_{n})}18\varepsilon_{0}n^{2}\cdot\eta_{ij}.
\end{eqnarray*}
If $\varepsilon_0 \le \frac{c_1p^2}{8}$, then by Lemma \ref{lem:badgood}, we thus have
\[
\sum_{ij\in E_{b}}\sum_{k\ell\in E_{g}}\eta(ij,k\ell)\le\frac{24\varepsilon_{0}}{c_{1}p^{2}}n^{2}\sum_{ij\in E_{g}}\eta_{ij}+\sum_{ij\in E(K_{n})}18\varepsilon_{0}n^{2}\cdot\eta_{ij}\le\frac{42\varepsilon_{0}}{c_{1}p^{2}}n^{2}\sum_{ij\in E(K_{n})}\eta_{ij}.
\]
Hence
\[
	\frac{42\varepsilon_{0}}{c_{1}p^{2}}n^{2}\sum_{ij\in E(K_{n})}\eta_{ij}
	\ge \sum_{ij\in E_{b}}\sum_{k\ell\in E_{g}}\eta(ij,k\ell) 
	> \frac{\beta}{1024c_{0}}n^{2}p\cdot\sum_{ij\in E_{b}}|\delta_{ij}|\|t_{ij}^{(0)}\|_2.
\]
If $\varepsilon_{0}\le\frac{\beta c_{1}^{2}p^{4}}{16 \cdot 42 \cdot 1024 c_{0}}$, then
\[
	\sum_{ij\in E(K_{n})}\eta_{ij}
	> \frac{\beta c_{1}p^{3}}{42\cdot1024c_{0}\varepsilon_{0}}\sum_{ij\in E_{b}}|\delta_{ij}|\|t_{ij}^{(0)}\|_2
	> \frac{16}{c_{1}p}\sum_{ij\in E_{b}}|\delta_{ij}|\|t_{ij}^{(0)}\|_2.
\]
If $\varepsilon_0 \le \frac{c_1p^2}{8}$, then by Lemma \ref{lem:transference}, this gives 
\[
	\sum_{ij\in E_{g}}\eta_{ij}\ge\frac{c_{1}p}{8}\sum_{ij\in E(K_{n})}\eta_{ij} 
	> 2\sum_{ij\in E_{b}}|\delta_{ij}|\|t_{ij}^{(0)}\|_2.
\]
Since Lemma \ref{lem:badgood} implies $\sum_{ij\in E_{g}}\eta_{ij}\ge 2 \sum_{ij\in E_{b}}\eta_{ij}$ 
(given $\eps_0 \leq \frac{c_1 p^2}{16}$), 
together with the inequality above, we have $\sum_{ij\in E_{g}}\eta_{ij}>\sum_{ij\in E_{b}}(|\delta_{ij}|\|t_{ij}^{(0)}\|_2+\eta_{ij})$.
By (\ref{eq:gain3-3d}), this shows that $R(T)>R(T_{0})$. 
The parameters must satisfy $\varepsilon_0 \le \min \{\frac{c_1p^2}{8}, \frac{\beta c_1^2 p^4}{16 \cdot 42 \cdot 1024 c_0}, \frac{c_1p^2}{16} \}$ and $\varepsilon_1 \le \frac{p}{192c_0}$.

\medskip{}


\noindent\textbf{Case 2}. $\overline{\delta} \neq 0$ and $\sum_{ij\in E_{g}'}|\delta_{ij}|\ge\frac{1}{8}\overline{\delta}|E_{g}'|$ and $|\Eb| \neq 0$.

In this case, we first show that there are large amount of positive
and negative parallel motions on the good edges. This will imply that
there is a large amount of rotational motions on the $K_{4}$s of
the form $\{i,j,k,\ell\}$ where $ij,k\ell\in E_{g}$ and $\delta_{ij}\ge0$,
$\delta_{k\ell}<0$. Since $\|t_{ij}^{(0)}\|_2 \ge \frac{1}{2}\mu$ for all $ij \in E_g$, 
Case 2 implies 
\[
	\sum_{ij\in E_{g}}|\delta_{ij}|\|t_{ij}^{(0)}\|_2
	\ge\sum_{ij\in E_{g}'}|\delta_{ij}|\|t_{ij}^{(0)}\|_2
	\ge\frac{1}{2}\mu\cdot\sum_{ij\in E_{g}'}|\delta_{ij}|
	\ge\frac{1}{16}\mu\overline{\delta}|E_{g}'|.
\]
Define $E_{+}=\{ij\in E_{g}\,:\,\delta_{ij}\ge0\}$ and $E_{-}=\{ij\in E_{g}\,:\,\delta_{ij}<0\}$.
The inequality above and (\ref{eq:good_parallel}) implies
\begin{align*}
	\sum_{ij\in E_{+}}\delta_{ij}\|t_{ij}^{(0)}\|_2 
	&= \frac{1}{2} \sum_{ij \in \Eg} (|\delta_{ij}| + \delta_{ij}) \| \toij\| \\
	&\ge \frac{1}{2}\left(\frac{1}{16}\mu\overline{\delta}|E_{g}'| - 2\sum_{ij \in E_b} |\delta_{ij}|\|\toij\|_2 \right)
  \ge  \frac{1}{2}\left(\frac{1}{16}\mu\overline{\delta}|E_{g}'|-2c_{0}\mu\overline{\delta}|E_{b}|\right).
\end{align*}
From \eqref{eq:longgoodedges}, we have $|E_{g'}| \ge \frac{1}{16c_0}n^2p$.
Therefore if $\varepsilon_0 \le \frac{p}{1024c_0^2}$, then 
\[
	\sum_{ij\in E_{+}}\delta_{ij}\|t_{ij}^{(0)}\|_2 
	\ge \frac{1}{32}\mu\overline{\delta}\left(\frac{1}{16c_0}n^{2}p-32c_0 \eps_0 n^2\right)
	\ge\frac{1}{1024c_0}\mu\overline{\delta}n^{2}p.
\]
Similarly $\sum_{ij\in E_{-}}(-\delta_{ij})\|t_{ij}^{(0)}\|_2\ge\frac{1}{1024c_0}\mu\overline{\delta}n^{2}p$.

We either have $|E_{+}|\ge\frac{1}{2}|E_{g}|$ or $|E_{-}|>\frac{1}{2}|E_{g}|$.
If the former holds, then by Lemmas \ref{lem:rigidity1} and \ref{lem:rigidity2},
\begin{eqnarray*}
\sum_{ij\in E_{-}}\sum_{k\ell\in E_{+}}\eta(ij,k\ell) & \ge & \sum_{ij\in E_{-}}\sum_{k\ell\in E_{+}\setminus B(ij)}\frac{\beta}{4}(-\delta_{ij})\|t_{ij}^{(0)}\|_2\\
 & \ge & \frac{\beta}{4}\cdot\sum_{ij\in E_{-}}(-\delta_{ij})\|t_{ij}^{(0)}\|_2(|E_{+}|-|B(ij)|).
\end{eqnarray*}
By Lemma \ref{lem:badK4}, we have $|B(ij)|\le6\varepsilon_{1}n^{2}$,
and thus $|E_{+}|-|B(ij)|\ge\frac{1}{2}|E_{g}|-6\varepsilon_{1}n^{2}\ge\frac{1}{2}(\frac{1}{4}n^{2}p-\varepsilon_{0}n^{2})-6\varepsilon_{1}n^{2}$.
If $\varepsilon_0<\frac{1}{16}p$ and $\varepsilon_{1}\le\frac{1}{192}p$, then
$|E_{+}|-|B(ij)|\ge\frac{1}{16}n^{2}p$, and the above gives
\[
	\sum_{ij\in E_{-}}\sum_{k\ell\in E_{+}}\eta(ij,k\ell)
	\ge\frac{\beta}{4}\cdot\frac{1}{16}n^{2}p\cdot\sum_{ij\in E_{-}}(-\delta_{ij})\|t_{ij}^{(0)}\|_2
	\ge\frac{\beta}{64}\cdot\frac{1}{1024c_0}\mu\overline{\delta}n^{4}p^{2}.
\]
Similarly, if $|E_{-}|>\frac{1}{2}|E_{g}|$, then $\sum_{ij\in E_{+}}\sum_{k\ell\in E_{-}}\eta(ij,k\ell)\ge\frac{\beta}{64\cdot1024c_0}\mu\overline{\delta}n^{4}p^{2}$.

On the other hand since each edge is contained in at most $\frac{n(n-1)}{2}$
copies of $K_{4}$ and $n$ copies of $K_{3}$ (and there are 6 edges
in a $K_{4}$), we have 
\[
\sum_{ij\in E_{-}}\sum_{k\ell\in E_{+}}\eta(ij,k\ell)\le\left(6\frac{n(n-1)}{2}+3n\right)\sum_{ij\in E(K_{n})}\eta_{ij}\le3n^{2}\sum_{ij\in E(K_{n})}\eta_{ij}.
\]
If $\varepsilon_{0}\le\frac{\beta c_{1}p^{3}}{32 \cdot 3 \cdot 64 \cdot 1024 \cdot c_{0}^2}$, then
\begin{eqnarray*}
	\sum_{ij\in E(K_{n})}\eta_{ij} 
	& \ge & \frac{1}{3n^{2}}\cdot\frac{\beta}{64\cdot1024c_0}\mu\overline{\delta}n^{4}p^{2}
	=\frac{\beta p^{2}}{3\cdot64\cdot1024c_0}\mu\overline{\delta}n^{2} \\
	& > & \frac{32}{c_{1}p}c_{0}\mu\overline{\delta}|E_{b}|
	\ge\frac{32}{c_{1}p}\sum_{ij\in E_{b}}|\delta_{ij}|\|t_{ij}^{(0)}\|_2,
\end{eqnarray*}
where the last inequality follows from Condition \ref{enu:length}.
If $\varepsilon_0 \le \frac{c_1p^2}{8}$, then by Lemma \ref{lem:transference}, this implies 
\[
	\sum_{ij\in E_{g}}\eta_{ij}
	\ge\frac{c_{1}p}{16}\sum_{ij\in E(K_{n})}\eta_{ij}
	>2\sum_{ij\in E_{b}}|\delta_{ij}|\|t_{ij}^{(0)}\|_2,
\]
Therefore from (\ref{eq:gain3-3d}) and (\ref{eq:badtwist-3d}), 
\begin{eqnarray*}
	R(T) & \ge & R(T_{0})+\sum_{ij\in E_{g}}\eta_{ij}-\sum_{ij\in E_{b}}(|\delta_{ij}|\|t_{ij}^{(0)}\|_2+\eta_{ij})\\
	& > & R(T_{0})+2\sum_{ij\in E_{b}}|\delta_{ij}|\|t_{ij}^{(0)}\|_2-\sum_{ij\in E_{b}}2|\delta_{ij}|\|t_{ij}^{(0)}\|_2=R(T_{0}).
\end{eqnarray*}
The parameters must satisfy $\varepsilon_0 \le \min\{\frac{p}{1024c_0^2}, \frac{1}{16}p, \frac{\beta c_1 p^3}{32 \cdot 3 \cdot 64 \cdot 1024 \cdot c_0^2}, \frac{c_1p^2}{8}\}$ and $\varepsilon_1 \le \frac{1}{192}p$.
\end{proof}



\subsection{Properties of Gaussians in three dimensions} \label{sec-gaussians-well-distributed-3d}

The first lemma establishes a bound on the average distance between random
Gaussian vectors.

\begin{lemma} \label{lem:gaussian_length}
%There exists a positive constant $c$ such that the following holds.
%\begin{itemize}
%  \setlength{\itemsep}{1pt} \setlength{\parskip}{0pt}
%  \setlength{\parsep}{0pt}
%\item[(i)] If $t_1, \ldots t_n \sim \mathcal{N}(0, I_{3 \times 3})$ are independent, then with probability at least $1 - 3e^{-cn}$
%\[
%	n\le\sum_{i\in[n]}\|t_{i}\|_2\le 3n.
%\]
%\item[(ii)] 
There exists a positive constant $c$ such that if $G$ is a $p$-typical graph with vertex set $[n]$, then with probability at least 
$1 - 3ne^{-cnp/2}$, 
\[
	\sum_{ij \in E(G)} \|t_i - t_j\|_2 \ge \frac{1}{8} n^2p.
\]
\end{lemma}
\begin{proof}
%(i)
%Each $\|t_i\|_2$ is subgaussian with mean $\sqrt{8/\pi}$.  By Proposition 5.10 in Vershynin\cite{Vershynin} on the concentration of subgaussians, there is a constant $c$ such that
%\[
%\P \Biggl( \Biggl| \sum_{i \in [n]} \|t_i\|_2 - n \sqrt{\frac{8}{\pi}} \Biggr| \geq \frac{n}{2} \Biggr) \leq e^{1-cn}
%\]
%Hence (i) follows.
%
%\medskip
%
%\noindent (ii)
Let $v \in \mathbb{R}^3$ be a fixed vector. 
Note that for all $j \in [n]$, we have
\[
	\|v - t_j\|_2^2
	= \|v\|^2 + \|t_j\|^2 - 2 \langle v, t_j \rangle.
\]
Therefore if $\langle v, t_j \rangle \le 0$, then $\|v - t_j\|_2 \ge \|t_j\|_2$. 
Further, by the symmetry of Gaussian random variables, we know that
the distribution of $\|t_j\|_2$ remains the same even after conditioning on the event
$\langle v, t_j \rangle \le 0$. Therefore
\[
	\mathbb{E}[\|v - t_j\|_2]
	\ge \P\Big(\langle v,t_j\rangle \le 0\Big) \cdot \mathbb{E}\Big[\|t_j\|_2 \,\Big|\, \langle v, t_j \rangle \le 0\Big]
	= \frac{1}{2} \mathbb{E}[\|t_j\|_2]
	= \sqrt{\frac{2}{\pi}},
\]
where the final equality holds since each $\|t_j\|_2$ is subgaussian with mean $\sqrt{8/\pi}$.  
Fix an index $i \in [n]$ and let $N_i$ be the neighborhood of $i$ in $G$. 
Since $G$ is $p$-typical, we have $|N_i| \ge \frac{1}{2}np$. 
By the analysis above, we see that
\[
	\mathbb{E}\left[ \sum_{j \in N_i} \|t_i - t_j\|_2 \right]
	\ge |N_i|\sqrt{\frac{2}{\pi}}.
\]
By Proposition 5.10 in Vershynin \cite{Vershynin} on the concentration of subgaussians, there is a constant $c$ such that
with probability at least $1 - e^{1-c|N_i|}$, we have
$\sum_{j \in N_i} \|t_i - t_j\|_2 \ge \frac{1}{2}|N_i| \ge \frac{1}{4}np$.
Therefore by taking the union bound over all indices $i \in [n]$, 
we see that with probability at least $1 - 3ne^{-cnp/2}$, 
\[
	\sum_{ij \in E(G)} \|t_i - t_j\|_2
	\ge \frac{1}{2} \sum_{i \in [n]}\sum_{j \in N_i} \|t_i - t_j\|_2
	\ge \frac{n}{2} \cdot \frac{1}{4}np = \frac{1}{8}n^2p. \qedhere
\]
\end{proof}

The second lemma establishes a bound on the angle between random Gaussian vectors.

\begin{lemma} \label{lem:gaussian_angle}
Let $x,y\in\mathbb{R}^{3}$ be linearly independent vectors.
If $t_{1},t_{2},\cdots,t_{n}\in\mathbb{R}^{3}$ are independent random
Gaussian vecotrs, then with probability $1-e^{-\Omega(\beta n)}$,
for all but at most $\beta n$ vectors $t_{i}$, we have 
\[
1-\left\langle \frac{t_{i}-x}{\|t_{i}-x\|_2},\,\frac{y-x}{\|y-x\|_2}\right\rangle ^{2}\ge\frac{\beta^{2}}{2(\|t_{i}\|_2^{2}+\|x\|_2^{2})}.
\]
\end{lemma}
\begin{proof}
Fix an index $i\in[n]$. Note that 
\begin{eqnarray}
1-\left\langle \frac{t_{i}-x}{\|t_{i}-x\|_2},\,\frac{y-x}{\|y-x\|_2}\right\rangle ^{2} & = & \left\Vert \proj_{(y-x)^{\perp}}\frac{t_{i}-x}{\|t_{i}-x\|_2}\right\Vert_2^{2}=\frac{\|\proj_{(y-x)^{\perp}}(t_{i}-x)\|_2^{2}}{\|t_{i}-x\|_2^{2}} \nonumber\\
 & \ge & \frac{\|\proj_{\{x,y\}^{\perp}}t_{i}\|_2^{2}}{\|t_{i}-x\|_2^{2}}\ge\frac{\|\proj_{\{x,y\}^{\perp}}t_{i}\|_2^{2}}{2(\|t_{i}\|_2^{2}+\|x\|_2^{2})}. \label{eq:ga_eq1}
\end{eqnarray}
Since $t_{i}$ is a random Gaussian vector and $x,y$ are linearly independent, 
the distribution of $\|\proj_{\{x,y\}^{\perp}}t_{i}\|_2$
is that of the absolute value of a standard normal distribution. Therefore
$\P(\|\proj_{\{x,y\}^{\perp}}t_{i}\|_2<\beta)\le\frac{1}{\sqrt{2\pi}}\int_{-\beta}^{\beta}e^{-x^{2}/2}dx\le\sqrt{\frac{2}{\pi}}\beta$.
Let ${\bf 1}_i$ be the indicator random variable of the event that
$\|\proj_{\{x,y\}^{\perp}}t_{i}\|_2<\beta$. We seen above that $\mathbb{E}[{\bf 1}_i] < \sqrt{\frac{2}{\pi}}\beta$.
Further, since $\{t_{i}\}_{i \in [n]}$ are independent, it follows that $\{{\bf 1}_i\}_{i \in [n]}$ are independent.
Therefore by Chernoff's inequality,
\[
	\P\left( \sum_{i \in [n]} {\bf 1}_i - \sqrt{\frac{2}{\pi}}\beta n > \frac{1}{5} \beta n \right)
	\le e^{-\Omega(\beta n)}.
\]
Hence with probability $1-e^{-\Omega(\beta n)}$,
there are at most $\sqrt{\frac{2}{\pi}}\beta n + \frac{1}{5} \beta n < \beta n$ vectors $t_{i}$ for which $\|\proj_{\{x,y\}^{\perp}}t_{i}\|_2<\beta$. The lemma now follows from \eqref{eq:ga_eq1}.
\end{proof}

The next lemma shows that random Gaussian vectors are well-distributed
with respect to a fixed pair of vectors.

\begin{lemma} \label{lem:gaussian_wd_prim}
There exists a positive real number $c$
such that the following holds for all pairs of linearly independent vectors $x,y\in\mathbb{R}^{3}$.
If $t_{1},\cdots,t_{n}\in\mathbb{R}^{3}$ are independent random Gaussian
vectors, then with probability $1-e^{-\Omega(n)}$,
the set of vectors $\{t_{1},\cdots,t_{n}\}$ are $\frac{c}{\max\{1,\|x+y\|_2\}}$-well-distributed
with respect to $(x,y)$. 
\end{lemma}
\begin{proof}
Let $c$ be a positive real number to be chosen later. We may rotate
the vectors so that $x=(\ell,0,x_{3})$ and $y=(\ell,0,y_{3})$ for
some $x_{3},y_{3},\ell\in\mathbb{R}$ where $\ell\ge0$. Note that
$\|x+y\|_2\ge2\ell$. Define $\ell_{0}=\max\{1,\ell\}$. It suffices
to give an estimate on the probability that
\[
\sum_{i=1}^{n}\|\proj_{span\{t_{i}-x,t_{i}-y\}^{\perp}}(h)\|_2\ge\frac{cn}{\ell_{0}}
\]
holds for all vectors $h\in(x-y)^{\perp}=\{(a,b,0):a,b\in\mathbb{R}\}$
satisfying $\|h\|_2=1$. 

Fix a vector $h=(h_{1},h_{2},0)$ satisfying $\|h\|_2=1$. For $t_{i}=(t_{i,1},t_{i,2},t_{i,3})$,
we have $span\{t_{i}-x,t_{i}-y\}=span\{(0,0,1),x+y-2t_{i}\}=span\{(0,0,1),(2\ell-2t_{i,1},-2t_{i,2},0)\}$.
Hence $s=(t_{i,2},\ell-t_{i,1},0)\in span\{t_{i}-x,t_{i}-y\}^{\perp}$, and
\begin{eqnarray*}
	\|\proj_{span\{t_{i}-x,t_{i}-y\}^{\perp}}(h)\|_2 
	& = & \frac{|\langle s,h\rangle |}{\|s\|_2}=\left|\frac{(t_{i,2},\ell-t_{i,1},0)\cdot h}{\sqrt{t_{i,2}^{2}+(\ell-t_{i,1})^{2}}}\right|\\
 & = & \frac{\left|h_{1}t_{i,2}+(\ell-t_{i,1})h_{2}\right|}{\sqrt{t_{i,2}^{2}+(\ell-t_{i,1})^{2}}}.
\end{eqnarray*}
Assume that $h_{1}\ge h_{2}\ge0$, which implies $h_{1}\ge\frac{1}{\sqrt{2}}$.
Since $t_{i,1}$ is normally distributed with variance 1, the probability
that $-1\le t_{i,1}\le0$ is $p$ for some fixed postive real number
$p$. Conditioned on this event and the event that $t_{i,2}\ge0$
(note that $t_{i,1}$ and $t_{i,2}$ are independent), we have 
\begin{eqnarray*}
	\|\proj_{span\{t_{i}-x,t_{i}-y\}^{\perp}}(h)\|_2
	& = & \frac{\left|h_{1}t_{i,2}+(\ell-t_{i,1})h_{2}\right|}{\sqrt{t_{i,2}^{2}+(\ell-t_{i,1})^{2}}}\ge\frac{h_{1}t_{i,2}}{\sqrt{t_{i,2}^{2}+4\ell_{0}^{2}}}.
\end{eqnarray*}
Therefore 
\[
	\P\left(\|\proj_{span\{t_{i}-x,t_{i}-y\}^{\perp}}(h)\|_2
	>\frac{1}{2\ell_{0}}\right)\ge\P\left(\frac{h_{1}t_{i,2}}{\sqrt{t_{i,2}^{2}+4\ell_{0}^{2}}}>\frac{1}{2\ell_{0}}\,\Big|\, t_{i,2}\ge0\right)\cdot\frac{1}{2}p.
\]
Note that for $t_{i,2}\ge0$, the inequality $\frac{h_{1}t_{i,2}}{\sqrt{t_{i,2}^{2}+4\ell_{0}^{2}}}>\frac{1}{2\ell_{0}}$
is equivalent to $h_{1}^{2}t_{i,2}^{2}>\frac{t_{i,2}^{2}}{4\ell_{0}^{2}}+1$,
which is equivalent to $t_{i,2}^{2}(h_{1}^{2}-\frac{1}{4\ell_{0}^{2}})>1$.
Since $h_{1}^{2}\ge\frac{1}{2}$ and $\ell_{0}\ge1$, we have 
\[
	\P\left(\|\proj_{span\{t_{i}-x,t_{i}-y\}^{\perp}}(h)\|_2>\frac{1}{2\ell_{0}}\right)
	\ge\P(t_{i,2}^{2}>4\,|\, t_{i,2}\ge0)\cdot\frac{1}{2}p=q
\]
for some fixed positive real number $q$. 
By considering the indicator random variable of the events 
$\|\proj_{span\{t_{i}-x,t_{i}-y\}^{\perp}}(h)\|_2>\frac{1}{2\ell_{0}}$, we see by Chernoff's inequality
that with probability $1-e^{-\Omega(n)}$, there are at least $\frac{qn}{2}$ indices 
$i \in [n]$ such that $\|\proj_{span\{t_{i}-x,t_{i}-y\}^{\perp}}(h)\|_2>\frac{1}{2\ell_{0}}$.
Note that this implies
\begin{eqnarray*}
	\sum_{i=1}^{n}\|\proj_{span\{t_{i}-x,t_{i}-y\}^{\perp}}(h)\|_2 
	%\ge \sum_{i=1}^{n}\P\left(\|\proj_{span\{t_{i}-x,t_{i}-y\}^{\perp}}(h)\|_2>\frac{1}{2\ell_{0}}\right)\cdot\frac{1}{2\ell_{0}} 
	\ge \frac{qn}{2} \cdot \frac{1}{2\ell_0}
	= \frac{qn}{4\ell_{0}}.
\end{eqnarray*}
To handle the case of $h_2 \geq h_1 \geq 0$, note that if $t_{i,1} \le 0$ and $0 \le t_{i,2} \le 1$, then
\begin{eqnarray*}
	\|\proj_{span\{t_{i}-x,t_{i}-y\}^{\perp}}(h)\|_2
	& = & \frac{\left|h_{1}t_{i,2}+(\ell-t_{i,1})h_{2}\right|}{\sqrt{t_{i,2}^{2}+(\ell-t_{i,1})^{2}}}
	\ge \frac{(\ell-t_{i,1})h_{2}}{\sqrt{1+(\ell-t_{i,1})^{2}}}
	\ge \frac{1}{\sqrt{2}} \cdot \frac{\ell-t_{i,1}}{\sqrt{1+(\ell-t_{i,1})^{2}}}.
\end{eqnarray*}
Since $\frac{x}{\sqrt{1+x^2}}$ is decreasing in the range $x \ge 0$, 
if $t_{i,1} \le -1$, then $\|\proj_{span\{t_{i}-x,t_{i}-y\}^{\perp}}(h)\|_2 \ge \frac{1}{2}$. 
Therefore we see as in above that with probability $1-e^{-\Omega(n)}$, there are at least $\frac{qn}{2}$ indices 
$i \in [n]$ such that $\|\proj_{span\{t_{i}-x,t_{i}-y\}^{\perp}}(h)\|_2>\frac{1}{2} \ge \frac{1}{2\ell_0}$.
All the remaining cases can be handled analogously.

Let $H$ be a set of $\lceil 2\pi \cdot \frac{8\ell_{0}}{q} \rceil \le \frac{60\ell_0}{q}$
vectors uniformly distributed along the circle $S_2 = \{(x,y,0) \,:\, x^2 + y^2 = 1\}$. 
Apply the analysis above to
each vector in $H$ and take the union bound to conclude that with
probability $1-\ell_{0}e^{-\Omega(n)}$, for all $h\in H$,
\[
\sum_{i=1}^{n}\|\proj_{span\{t_{i}-x,t_{i}-y\}^{\perp}}(h)\|_2\ge\frac{q}{4}\frac{n}{\ell_{0}}.
\]
Let $h' \in S_2$ be an arbitrary vector and
let $h \in H$ be the vector closest to $h'$. The distance from $h$ to $h'$ along the circle $S_2$ is at most 
$2\pi \cdot \frac{1}{|H|} \le \frac{q}{8\ell_0}$, and hence $\|h - h'\|_2 \le \frac{q}{8\ell_0}$. 
Thus for all $i$, we have 
\begin{eqnarray*}
\|\proj_{span\{t_{i}-x,t_{i}-y\}^{\perp}}(h')\|_2 & \ge & \|\proj_{span\{t_{i}-x,t_{i}-y\}^{\perp}}(h)\|_2-\|h-h'\|_2\\
 & \ge & \|\proj_{span\{t_{i}-x,t_{i}-y\}^{\perp}}(h)\|_2-\frac{q}{8\ell_{0}}
\end{eqnarray*}
Therefore 
\[
\sum_{i=1}^{n}\|\proj_{span\{t_{i}-x,t_{i}-y\}^{\perp}}(h')\|_2\ge\frac{q}{4}\frac{n}{\ell_{0}}-\frac{q}{8}\frac{n}{\ell_{0}}\ge\frac{q}{8}\frac{n}{\ell_{0}}.
\]
Since $\ell_{0}=\max\{1,\ell\}\le\max\{1,\|x+y\|_2\}$, this implies
the lemma for $c=\frac{q}{8}$.
\end{proof}


By applying the union bound together with the three lemmas above,
we obtain the following lemma.

\begin{lemma} \label{lem:gaussian_wd}
There exists $c, \zeta \in \mathbb{R}$ and $n_0 \in \mathbb{N}$ such
that the following holds for all positive real numbers $\varepsilon$ and natural numbers $n \ge n_0$.
Let $G$ be a $p$-typical graph with vertex set $[n]$ for some $p$ satisfying 
$np^2 \geq \zeta \log n$. If $t_{1},\cdots,t_{n}\in\mathbb{R}^{3}$
are independent random Gaussian vectors, then the following holds
with probability $1-n^{-5}$,
\begin{itemize}
  \setlength{\itemsep}{1pt} \setlength{\parskip}{0pt}
  \setlength{\parsep}{0pt}
\item[1.] For each distinct $i,j \in [n]$, for all but at most $\varepsilon n$ indices $k \in [n]$, we have $1 - \langle \frac{t_k - t_i}{\|t_k - t_i \|_2}, \frac{t_j - t_i}{\|t_j - t_i\|_2} \rangle^2 \ge \frac{\varepsilon^2}{64 \log n}$,
\item[2.] for all distinct $i,j \in [n]$, we have $\|t_i - t_j\|_2 \le 40\sqrt{\log n} \cdot \mu$, where $\mu = \frac{1}{|E(G)|} \sum_{ij \in E(G)} \|t_i - t_j\|_2$, and
\item[3.] the set $\{t_i\}_{i\in[n]}$ is $\frac{c}{\sqrt{\log n}}$-well-distributed along $G$.
\end{itemize}\end{lemma}
\begin{proof}
Let $c$ be eight times the constant coming from Lemma \ref{lem:gaussian_wd_prim}.
For each distinct $i,j\in[n]$, define $S_{ij}=\{t_{k}\,:\, ik,jk\in E(G)\}$.
Consider the following events:
\begin{itemize}
  \setlength{\itemsep}{1pt} \setlength{\parskip}{0pt}
  \setlength{\parsep}{0pt}
\item[(i)] for all $i\in[n]$, we have $\|t_i\|_2 \le  4\sqrt{\log n}$,
\item[(ii)] $\sum_{ij \in E(G)} \|t_i - t_j \|_2 \ge \frac{1}{8}n^2p$, %$n \le \sum_{i\in[n]} \|t_i\|_2 \le 3n$ and 
\item[(iii)] for each distinct $i,j \in [n]$, for all but at most $\varepsilon n$ integers $k \in [n]$, we have $1 - \langle \frac{t_k - t_i}{\|t_k - t_i \|_2}, \frac{t_j - t_i}{\|t_j - t_i\|_2} \rangle^2 \ge \frac{\varepsilon^2}{2(\|t_k\|_2^2 + \|t_i\|_2^2)}$,
\item[(iv)] for each distinct $i,j\in[n]$, $S_{ij}$ is $\frac{8c}{\max\{1,\|t_i+t_j\|_2\}}$-well-distributed with respect to $(t_i,t_j)$.
\end{itemize}
For a fixed $i \in [n]$, since $\|t_i\|_2^2$ follows a $\chi^2$ distribution with $3$ degrees of freedom, 
standard estimates on Chi-squared random variables, such as Lemma 1 in \cite{LaurentMassart}, give
\[
\P( \|t_i\|_2^2 \geq 3 + 2 \sqrt{3}t + 2t^2 ) \leq e^{-t^2}.
\]
Let $t = \sqrt{7 \log n}$. 
If $n$ is sufficiently large, then 
$2t^2 + 2 \sqrt{3} t + 3 < 16 \log n$. 
As a result, $\P(\|t_i\|_2^2 \geq 16 \log n) \leq e^{-7 \log n}$.
Hence
Property (i) holds with probability $1-n^{-6}$ by taking the union
bound over all $i \in [n]$. 
Property (ii) holds with probability $1-e^{-\Omega(n)}$ by
Lemma \ref{lem:gaussian_length}. For a fixed pair $i,j\in[n]$, Property
(iii) holds with probability $1-e^{-\Omega(\varepsilon n)}$ by Lemma \ref{lem:gaussian_angle}.
Hence by taking the union bound, we see that Property (iii) holds
with probability $1-n^{2}e^{-\Omega(\varepsilon n)}$. For a fixed pair $i,j\in[n]$,
by Lemma~\ref{lem:gaussian_wd_prim}
and the fact that each pair is contained in at least $\frac{1}{2}np^2$ triangles, 
we have Property (iv) for the pair $i,j$
with probability $1-e^{-\Omega(np^2)}$. Hence by taking
the union bound, we see that Property (iv) holds with probability
$1-n^{2}e^{-\Omega(np^2)}$. 
Thus we see that all four
events (i)-(iv) simultaneously hold with at least probability $1-n^{-5}$ for sufficiently large $n$, provided that $np^2 \geq \zeta \log n$ for sufficiently large $\zeta$.

We now show that Properties (i)-(iv) imply Properties 1-3. Note that
Properties 1 and 3 immediately follow from Properties (i), (iii), and
(iv). Further, since $|E(G)|\le n^2p$, Property (ii) implies
\begin{eqnarray*}
	\mu & = & \frac{1}{|E(G)|}\sum_{ij\in E(G)}\|t_{i}-t_{j}\|_2
	\ge\frac{1}{n^2p} \cdot \frac{1}{8}n^2 p = \frac{1}{8}.
\end{eqnarray*}
Hence by Property (i), we have for all $i,j\in[n]$, 
\[
	\|t_{i}-t_{j}\|_2
	\le \|t_{i}\|_2+\|t_{j}\|_2
	\le 8\sqrt{\log n}\le 64\mu \sqrt{\log n}.\qedhere
\]
\end{proof}




\subsection{Proof of Theorem \ref{thm-3d}} \label{sec-proof-random-theorem-3d}

We can now prove the three-dimensional recovery theorem, which we state here again for convenience:
\setcounter{theorem}{1}
\begin{theorem}
There exists $n_0 \in \mathbb{N}$ and $c \in \mathbb{R}$ such that the following holds for all $n \ge n_0$.
Let $G([n],E)$ be drawn from $G(n,p)$ for some $p = \Omega( n^{-1/5} \log^{3/5} n)$. Take $ \tnot_1, \ldots \tnot_n  \in \R^3$, where $\toi \sim \mathcal{N}(0, I_{3 \times 3})$ are i.i.d., independent from $G$. 
%If $n_0 \leq  n  $, then 
There exists $\gamma = \Omega(p^5 / \log^3 n)$ and an event of probability at least  $1- \frac{1}{n^{4}}$ %, where $c>0$ is an absolute constant, 
on which the following holds:\\[1em]
For arbitrary  subgraphs $\Eb$ satisfying $\max_i \deg_b(i) \leq \gamma n$ and arbitrary pairwise direction corruptions $\vij \in \mathbb{S}^2$ for $ij \in \Eb$,  the convex program \eqref{shapefit} has a unique minimizer equal to $\left \{\alpha \Bigl(\toi - \tnotbar \Bigr)\right\}_{i \in [n]}$ for some positive $\alpha$ and for $\tnotbar = \frac{1}{n}\sum_{i\in [n]} \toi$. 
\end{theorem}

%\begin{theorem}
%There exists a constant $c \in \mathbb{R}$ and $n_0 \in \mathbb{N}$
%such that the following holds. If the
%measurement graph is $G(n,p)$ and the input is a set of independent
%random Gaussian vectors in $\mathbb{R}^{3}$, then with probability
%at least $1-n^{-4}$, our algorithm successfully recovers when there
%are at most $\gamma \frac{p^{5}}{\log^{3}n}n$ bad edges incident
%to each vertex.
%\end{theorem}
\begin{proof}
Let $n_0$ be a sufficiently large natural number larger than
that coming from Lemma~\ref{lem:gaussian_wd}.
Lemma~\ref{lem:ptyp}  implies $G$ is $p$-typical with probability
$1 - n^2 e^{-\Omega(np^2)}$. Condition on $G$ being $p$-typical.
Let $c$ be the constant from Lemma \ref{lem:gaussian_wd}.
By applying Lemma \ref{lem:gaussian_wd}
with $\varepsilon=\frac{p}{2^{15}\sqrt{\log n}}$, with probability
at least $1-n^{-5}$, we have
\begin{itemize}
  \setlength{\itemsep}{1pt} \setlength{\parskip}{0pt}
  \setlength{\parsep}{0pt}
\item[1.] For each distinct $i,j \in [n]$ satisfying $i < j$, for all but at most $2 \eps n = \frac{p}{2^{14}\sqrt{\log n}}$ integers $k \in [n]$, we have $1 - \langle \frac{t_k - t_i}{\|t_k - t_i \|}, \frac{t_j - t_i}{\|t_j - t_i\|} \rangle^2 \ge \frac{p^2}{2^{36} \log^2 n}$ and $1 - \langle \frac{t_k - t_j}{\|t_k - t_j \|}, \frac{t_i - t_j}{\|t_i - t_j\|} \rangle^2 \ge \frac{p^2}{2^{36} \log^2 n}$,
\item[2.] for all distinct $i,j \in [n]$, we have $\|t_i - t_j\| \le 64\sqrt{\log n} \cdot \mu$, where $\mu = \frac{1}{|E(G)|} \sum_{ij \in E(G)} \|t_i - t_j\|$, and
\item[3.] the set $\{t_i\}_{i\in[n]}$ is $\frac{c}{\sqrt{\log n}}$-well-distributed along $G$.
\end{itemize}
Thus the probability that $G$ is $p$-typical and Properties 1-3 listed above holds
is at least $1 - n^{-4}$. 
Hence we may apply Theorem \ref{thm-deterministic-three-d} with with 
$c_{0}=64\sqrt{\log n}$,
$\varepsilon_{1}= \frac{p}{2^{14}\sqrt{\log n}} \le \frac{p}{192c_{0}}$,
$\beta=\sqrt{\frac{p^2}{2^{36}\log^{2}n}}=\frac{p}{2^{18}\log n}$,
and $c_{1}=\frac{c}{\sqrt{\log n}}$. The theorem holds if 
\[
	\varepsilon_{0}
	\le \frac{c^2p^5}{2^{53} \log^{3}n}
	\le  \frac{p}{2^{18}\log n} \cdot \frac{c^2}{\log n} p^4 \cdot \frac{1}{32 \cdot 3 \cdot 64 \cdot 1024 \cdot 64^2 \log n}
		= \frac{\beta c_1^2 p^4}{32 \cdot 3 \cdot 64 \cdot 1024\  c_0^2}.		
\]
Letting $\gamma$ from the theorem statement be $\eps_0$, note that the condition $\max_i \deg_b(i) \leq \gamma n$ is nontrivial when $p = \Omega(n^{-1/5} \log^{3/5} n)$.  \qedhere
\end{proof}







